%%
%% This is a sample doctoral dissertation.  It shows the appropriate
%% structure for your dissertation.  It should handle most of the
%% strange requirements imposed by the Grad school; like the different
%% handling of titles of one/many appendices.  It will automatically
%% handle the linespacing changes.  The body default is double-spaced
%% (except when you use the singlespace or condensed options).  The
%% default for quotations is single-space, and the default for tabular
%% environments is also single-space.
%%
%% This class adds the following commands and environments to the
%% report class, upon which it is based:
%% Commands
%% ------------
%% \degree{name}{abbrv} -- Sets the name and abbreviation for the degree.
%%                         These default to ``Doctor of Philosopy''
%%                         and ``Ph.D.'', respectively.
%% \copyrightyear{year} -- for the copyright page.
%% \bachelors{degree}{institution} -- for the abstract
%% \masters{degree}{institution}   --  "
%%     if you have other degrees you may use
%% \secondbachelors{degree}{institution}
%% \thirdbachelors{degree}{institution}
%% \secondmasters{degree}{institution}
%% \thirdmasters{degree}{institution}
%% \priordoctorate{degree}{institution}
%%
%% \committeechair{name}           -- for the signature page
%% or, if you have two co-chairs:
%% \cochairs{first name}{second name}
%%
%% \firstreader{name}              --  "
%% \secondreader{name}             --  "
%% \thirdreader{name}              -- (optional)
%% \fourthreader{name}             --  "
%% \fifthreader{name}              --  "
%% \sixthreader{name}              --  "
%% \departmentchair{name}          -- for the signature page
%% \departmentname{name}           --  "
%%
%% \copyrightpage                  -- produces the copyright page
%% \signaturepage                  -- produces the signature page
%%
%% \frontmatter                    -- these are required in their various
%% \mainmatter                     -- appropriate locations
%% \backmatter                     --
%%
%% \unnumberedchapter[toc]{name}   -- like \chapter, except that it
%%                                    produces an unnumbered chapter;
%%                                    alternatively, like \chapter*,
%%                                    except that it lists the chapter
%%                                    in the table of contents.
%%
%% New environments:
%%   dedication  -- for the dedication
%%   abstract    -- for the abstract
%%
%% The thesis documentclass is built on top of the report document class.
%% It accepts all of the options that the report class accepts, plus the
%% following:
%%     doublespace -- the default, indicates double spacing as per U.Mass.
%%                    requirements.  You will need this when you do your
%%                    final copy.
%%     singlespace -- for earlier work, not acceptable to the Grad school
%%     condensed   -- for earlier work, not acceptable to the Grad school,
%%                    creates condensed versions of the frontmatter.
%%                    Condensed implies singlespace.
%%     dissertation - the default, indicates that this document is a
%%                    dissertation.
%%     proposal    -- indicates that this document is a dissertation proposal,
%%                    rather than a dissertation.  This will only change the
%%                    wording on the title and signature pages.
%%     thesis      -- indicates that this document is a Master's thesis
%%                    rather than a doctoral dissertation.  This also changes
%%                    the default for \degree to Master of Science, M.S.
%%     allowlisthypenation -- (the default), allows hyphenation of words in
%%                    the table of contents, the list of figures, and the list
%%                    of tables.  I believe that this is acceptable to the
%%                    Graduate School.
%%     nolisthyphenation -- disallows hyphenation of words in the table of
%%                    contents and the list of figures and tables.  Use this
%%                    option if the Grad School doesn't like your hyphenation.
%%     nicerdraft  -- relaxes some of the Grad School's rules for working with
%%                    drafts -- has no effect when doublespace is in effect
%%     nonicerdraft -- the default, leaves things in draft as they will be in
%%                     the final version
%% umassthesis changes the default font size to 12pt, but you may specify 10pt or
%%   11pt in the options.
%% \documentclass{umassthesis}          % for Ph.D. dissertation or proposal
 \documentclass[thesis]{umassthesis}  % for Master's thesis

%%
%% If you have enough figures or tables that you run out of space for their
%% numbers in the List of Tables or List of figures, you can use the following
%% command to adjust the space left for numbers.  The default is shown:
%%
%% \setlength{\tablenumberwidth}{2.3em}

%% Use the hyperref package if you're producing a version for online
%% distribution and you want hyperlinks.  Note that the Grad School doesn't want
%% their PDF viewers to colorize or otherwise highlight the links; use the
%% hidelinks option to hyperref to avoid decorating links.
\usepackage[hidelinks]{hyperref}
\usepackage{graphicx}
\usepackage{amsmath}
\usepackage{amssymb}
\usepackage{lipsum}
\usepackage{amsmath}
\usepackage{latexsym}
% Useful for Knitr. Otherwise, I use memoir class's 'ctabular' env.
\usepackage{longtable}

%% One way of formatting the epigraph/frontispiece is to use this package.
%\usepackage{epigraph}

\usepackage{Sweave}
\begin{document}
\Sconcordance{concordance:ms_thesis.tex:ms_thesis.Rnw:%
1 116 1 1 0 3 1 1 20 191 1 2 2 99 1 1 7 8 0 1 2 91 1 1 4 33 0 1 2 6 1 1 %
4 33 0 1 2 9 1 2 2 12 1 1 4 16 0 1 2 6 1 1 4 16 0 1 2 3 1 1 5 11 1 2 2 %
12 1 2 2 12 1 2 2 5 1 2 2 8 1 2 2 47 1 1 5 36 0 1 2 3 1 1 6 5 1 1 5 36 %
0 1 2 3 1 1 6 6 1 1 6 49 0 1 2 3 1 1 7 5 1 1 6 49 0 1 2 3 1 1 7 3 1 1 5 %
3 1 1 6 89 0 1 2 3 1 1 8 11 1 1 6 89 0 1 2 4 1 1 5 20 1}



%%
%% You must fill in all of these appropriately
\title{Forecasting COVID-19 with\protect\\Temporal Hierarchies and Ensemble Methods}
\author{Li Shandross}
\date{May 2023} % The date you'll actually graduate -- must be
                     % February, May, or September
\copyrightyear{2023}
\bachelors{B.Sc.}{University of Massachusetts Amherst}
\committeechair{Nicholas G. Reich}
\firstreader{Evan L. Ray}
\secondreader{Benjamin W. Rogers}
%% \thirdreader{Bill Shepherd}
\departmentchair{Lisa Chasan-Taber} % Default uses "Department Chair" as the title. To
% use an alternate title, such as "Chair", use \departmentchair[Chair]{Pete Shearer}
% CICS uses "Chair of the Faculty" as of 2019.
\departmentname{Biostatistics and Epidemiology}
% \departmentname{Robert and Donna Manning College of \\Information and Computer Sciences}

%% If your degree is something other than a Ph.D. (for a dissertation), or
%% an M.S. (for a thesis), you will need to uncomment the appropriate
%% following line:
%%
%% \degree{Doctor of Education}{Ed.D.}
%% \degree{Doctor of Philosophy}{Ph.D.}
%%
%% \degree{Master of Arts}{M.A.}
%% \degree{Master of Arts in Teaching}{M.A.T.}
%% \degree{Master of Business Administration}{M.B.A.}
%% \degree{Master of Education}{M.Ed.}
%% \degree{Master of Fine Arts}{M.F.A.}
%% \degree{Master of Landscape Architecture}{M.L.A.}
%% \degree{Master of Music}{M.M.}
%% \degree{Master of Public Administration}{M.P.A.}
%% \degree{Master of Public Health}{M.P.H.}
%% \degree{Master of Regional Planning}{M.R.P.}
%% \degree{Master of Science}{M.S.}
%% \degree{Master of Science in Accounting}{M.S. Acctg.}
%% \degree{Master of Science in Chemical Engineering}{M.S. Ch.E.}
%% \degree{Master of Science in Civil Engineering}{M.S.C.E.}
%% \degree{Master of Science in Electrical and Computer Engineering}{M.S.E.C.E.}
%% \degree{Master of Science in Engineering Management}{M.S. Eng. Mgt.}
%% \degree{Master of Science in Environmental Engineering}{M.S. Env. E.}
%% \degree{Master of Science in Industrial Engineering and Operations Research}{M.S.I.E.O.R.}
%% \degree{Master of Science in Manufacturing Engineering}{M.S. Mfg. Eng.}
%% \degree{Master of Science in Mechanical Engineering}{M.S.M.E.}
%%
%% \degree{Professional Master of Business Administration}{P.M.B.A.}


%%
%% These lines produce the title, copyright, and  signature pages.
%% They are Mandatory; except that you could leave out the copyright page
%% if you were preparing an M.S. thesis instead of a PhD dissertation.
\frontmatter
\maketitle
%% \copyrightpage     %% not required for an M.S. thesis
\signaturepage

%%
%% Dedication is optional -- but this is how you create it
%% \begin{dedication}              % Dedication page
%%   \begin{center}
%%     \emph{To those little lost sheep.}
%%   \end{center}
%% \end{dedication}

%%
%% Epigraph (aka frontispiece) is also optional, but this is one way you
%% can create it
%\begin{frontispiece}
%  %% Format to your liking -- see documentation of epigraph package
%  \setlength{\epigraphrule}{0pt}
%
%  \begin{epigraphs}
%    \qitem{%
%      \itshape
%      Mary had a little lamb,\\
%      Her fleece was white as snow.\\
%      \vspace{\baselineskip}
%      And everywhere that Mary went\\
%      The lamb was sure to go.
%      \vspace{\baselineskip}}
%    {Sarah Josepha Hale}
%
%    \vspace{2\baselineskip}
%    \qitem{%
%      \itshape
%      Baa, baa, black sheep,\\
%      Have you any wool?\\
%      Yes, sir, yes, sir,\\
%      Three bags full;\\
%      One for the master,\\
%      And one for the dame,\\
%      And one for the little boy\\
%      Who lives down the lane.
%      \vspace{\baselineskip}}
%    {English Nursery Rhyme}
%
%  \end{epigraphs}
%\end{frontispiece}

%%
%% Acknowledgements are optional...yeah, right.
%% \chapter{Acknowledgments}             % Acknowledgements page
%%   Thanks to all those fine shepherds. Not to mention all the great
%%    border collies and suchlike fine animals.

%%
%% Abstract is MANDATORY. -- Except for MS theses
\begin{abstract}                % Abstract
Infectious disease forecasting efforts underwent rapid growth during the COVID-19 pandemic, providing guidance for pandemic response and about potential future trends. Yet despite their importance, short-term forecasting models often struggled to produce accurate real-time predictions of a complex and rapidly changing system. This gap in accuracy, which persists three years into the pandemic, warrants the exploration and testing of new methods to glean fresh insights.

In this work, we examine the application of the temporal hierarchical forecasting (THieF) methodology to probabilistic forecasts of COVID-19 incident hospital admissions in the United States. THieF is a recent, innovative forecasting technique that aggregates time-series data into a hierarchy made up of different temporal scales, produces forecasts at each level of the hierarchy, then reconciles those forecasts using optimized weighted forecast combination. While THieF's unique approach has shown substantial accuracy improvements in a diverse range of applications, such as operations management and emergency room admission predictions, this technique has not previously been applied to outbreak forecasting.

We generate candidate models formulated using the THieF methodology, which differ by their hierarchy schemes and data transformations, and weighted ensembles of the THieF models. The models are validated using weighted interval score (WIS), and the top-performing subset is compared to several benchmark models. These models include simple SARIMA models, a naive baseline model, and a median ensemble of operational incident hospitalization models from the US COVID-19 Forecast Hub. The THieF models and THieF ensembles demonstrate improvements in WIS and MAE, as well as competitive prediction interval coverage, over the benchmark models for both the validation and testing phases. These accuracy improvements suggest the THieF methodology may serve as a promising new technique for infectious disease forecasting that makes better predictions in real-time.
\end{abstract}

%%
%% Preface goes here...would be just like Acknowledgements -- optional
%% \chapter{Preface}
%% ...

%%
%% Table of contents is mandatory, lists of tables and figures are
%% mandatory if you have any tables or figures; must be in this order.
\tableofcontents                % Table of contents
\listoftables                   % List of Tables
\listoffigures                  % List of Figures

%%
%% We don't handle List of Abbreviations
%% We don't handle Glossary


%%%%%%%%%%%%%%%%%%%%%%%%%%%%%%%%%%%%%%%%%%%%%%%%%%%%%%%%%%%%%%%%%%%%%%%%%
%% Time for the body of the dissertation
\mainmatter   %% <-- This line is mandatory

\chapter{Introduction}

\section{Infectious disease forecasting}
\section{COVID-19 forecasting}
\subsection{US COVID-19 Forecast Hub}
\section{Temporal hierarchical forecasting}
\section{Goals}

\chapter{Methods}
\section{Surveillance data}
This analysis uses confirmed hospital admissions from HealthData.gov from the COVID-19 Reported Patient Impact and Hospital Capacity by State Time Series. This dataset reports state-aggregated daily incident hospitalizations starting on January 1, 2020 for U.S. states and territories. The data come from three primary sources: HHSTeleTracking, state/territorial health departments which report to HHS Protect, and National Health Care Safety Network (prior to July 15, 2020) \cite{hhs_data}.

On December 1, 2020, the U.S. COVID-19 Forecast Hub designated this dataset as the official ground truth data for COVID-19 incident hospitalizations. (No other sources were designated official ground truth data prior to this date.) \cite{hub_data_repo} This meant that the HealthData.gov data was used to evaluate the performance of real-time incident hospitalization forecasts submitted to the Forecast Hub. Hence, teams were encouraged to use this hospitalization data set to train their models \cite{cramer}. In the analysis, we also use this dataset as both input training data for our models and truth data against which to perform accuracy evaluations.

Hospitalization reporting was sporadic for the first four months of the pandemic, with consistent, reliable values for major states and territories achieved only in late July 2020. Specifically, reports of daily confirmed hospitalizations for all 50 states, Washington D.C., Puerto Rico, and the U.S. as a whole began on July 27, 2020. Although COVID-19 hospitalization reporting stabilized after that for cases and deaths, hospitalization values are often more reliable due to requirements for Medicare reimbursement. COVID-19 cases have been vastly underreported for the entire duration of the pandemic \cite{mech_bayes}, and this problem has worsened with dwindling testing. Further, truth data collection and reporting for COVID-19 incident cases and deaths has been discontinued by Johns Hopkins University (the official ground truth data source) as of March 10, 2023 \cite{jhu_covid}.

Surveillance data, like the HealthData.gov confirmed hospitalizations, sometimes required revisions due to reporting anomalies. Hospitalizations might be entered on incorrect dates or added after COVID-19 was determined to be the cause of hospitalization. Updates to values usually occurred within the first week or so of posting. As a result, real-time forecast models might train on data that include incident hospitalization numbers for location-date combinations that do not reflect the finalized values.

We query these confirmed hospitalizations from the `covidData` package through the `covidHubUtils` package, both developed by Reich Lab at the University of Massachusetts Amherst to deal with COVID-19 surveillance data. Versions of this data are stored such that users may access the truth data as of a specified date. For example, it is possible to retrieve of version of the confirmed hospitalization data set as of July 27, 2020. To ensure fair comparisons between the models developed for this analysis and some real-time models, our models are trained on the confirmed hospitalizations that would have been available as of the "forecast date" of the retrospective forecast.

\section{Forecasts}
A forecast for a target-location combination (e.g. "1-day ahead incident hospitalizations in Massachusetts") is defined by one or both of a point forecast and a probabilistic forecast, represented by a set of 23 quantiles \cite{cramer}. (More information is given about forecast targets, horizons, locations, and quantiles in the following subsections.) According to conventions set by the Forecasts Hub, all forecast values are truncated to be non-zero, as there is no real-world meaning for negative incident hospitalizations \cite{ensembles_info}. Further, while this is a retrospective analysis, forecasts are made on Sunday weekly (as they would be if made in real-time) using the data available as of that day. We chose to make forecasts on Sunday because confirmed hospitalization values for subsequent days did not update until Tuesday or later until some time in March 2021, which presented several issues. First, if we began validation forecasting after March 2021, we would lose three months of the period of analysis. Second, the main comparison models made forecasts on Monday, meaning that any models that forecast the week before would be add a disadvantage while any models that forecast on Tuesday or later would be evaluated with the following week's forecasts. Third, THieF does not allow for missing data points.

\subsection{Forecast horizons}
Short-term forecasts for COVID-19 have been defined to encapsulate horizons of 1- to 4-weeks ahead, or, equivalently, 1- to 28-days ahead \cite{cramer}. Forecasts for incident hospitalizations are made on a daily time scale, so we choose to restrict\footnote{Some models make up to 84-day ahead forecasts due to the THieF methodology construction.} the forecast horizons we evaluate to 1- to 28-days ahead.\footnote{Note, there is a slight exception to this restriction. See Metrics and Evaluation for more information.}

\subsection{Forecast locations}
We make forecasts for 53 of 55 locations from the HealthData.gov confirmed hospitalizations dataset during the period of analysis. These 55 locations include the 50 states, four jurisdictions/territories (Washington D.C., Puerto Rico, American Samoa, and Guam), and a U.S. national level. American Samoa and Guam are excluded due to their very low or zero COVID-19 incident hospitalizations over the entire pandemic. Forecasts for low-count locations that predicted single-digit or zero hospitalizations were often very accurate; hence, these locations do not offer meaningful contributions to our understanding of differences in forecast skill between modeling approaches \cite{cramer}. Further, even though we include the U.S. national level, the remaining 52 locations are our primary focus, as the magnitude of national forecasts is much greater than out of any single state.

\subsection{Forecast quantiles}
The Forecast Hub specifies the usage of 23 quantiles \[\mathbb{Q} = \{.010, 0.025, .050, .100, \cdots, .900, .950, .990\}\] to represent a full probabilistic distribution, with the median 0.500 quantile usually taken as the model's point forecast. We follow this construction, though with a slight modification: prediction intervals are used to define a full distribution of forecasts. The bounds of these intervals can be used to obtain the desired quantiles, as shown in the table below. For example, the values for a 0.500 and a 0.950 quantile (see the Quantile row) are obtained by taking the lower and upper bounds of the 95\% prediction interval (see the Interval row, where "L" and "U" represent the lower and upper bounds, respectively).

\begin{table}[h]
\begin{center}
\label{tab:table_quant_int}
% latex table generated in R 4.1.3 by xtable 1.8-4 package
% Tue Mar 21 18:58:31 2023
\begin{tabular}{rllllllllllll}
  \hline
 &  \\ 
  \hline
Quantile & .010 & .025 & .050 & 0.100 & 0.150 & ... & 0.500 & ... & 0.900 & 0.950 & .975 & .990 \\ 
  Interval & 98L & 95L & 90L & 80L & 70L & ... & Median & ... & 80U & 90U & 95U & 98U \\ 
   \hline
\end{tabular}\caption{Table of the 23 quantiles required by the Forecast Hub and their associated prediction interval bound.}
\end{center}
\end{table}

\section{Date range of analysis}
We focused on the period between Monday, July 27, 2020 and Sunday, October 2, 2022. This date range spans 798 days, or 114 weeks, which is just over 26 months. The entire period of analysis is split into two main phases: a validation phase of 462 days (66 weeks) and a testing phase of 236 days (48 weeks). The split between these two phases is such that the validation phase ends on Sunday, October 31, 2021 and the testing phase begins on Monday, November 1, 2021.

The first day of the period of analysis (July 27, 2020) was selected as the first day of reliable truth data for all locations of interests. Forecasting for the validation phase began on Monday, December 7, 2020, as this is the first Monday after the HealthData.gov confirmed hospitalizations data set was declared the official ground truth data. The division between the validation and testing phases was chosen to encapsulate three COVID-19 variant-based waves in the validation phase (Winter 2020-21, Alpha, and Delta) and two waves in the testing phase (Omicron and BA.4/BA.5). The separation between pandemic waves is made based on U.S. national level incident hospitalizations, pictured below in [Figure X].

\begin{figure}
\label{fig:figure_evaluate_period_plot"}
% \includegraphics[width=0.5\textwidth, inner]{overleaf-logo}
\includegraphics{ms_thesis-evaluation_period_plot}
\caption{Overview of the period of analysis included in the paper, superimposed over COVID-19 incident hospitalizations for the entire United States. Vertical dashed lines indicate divisions between phases. The training phase begins on Monday, July 27, 2020 and ends on Monday, July 5, 2021. This phase covers 343 days, or 49 weeks. Forecasts for the training phase are made starting on Monday, December 7, 2020. The testing phase, although not included in this paper, begins Tuesday, July 6, 2021 and ends Monday, November 1, 2021. The division is based on including two waves (the Winter 2020-21 and the Alpha variant) of the pandemic in the training phase and one (Delta variant) in the testing phase. July 27, 2020 marks the beginning of when the forecast locations of interest had reliable truth data.}
\end{figure}


\section{Benchmark comparison models}
Two models from the Forecast Hub---the COVIDhub-baseline and COVIDhub-4\_week\_ensemble---are selected to serve as comparisons for the models in this analysis during the testing phase evaluation. Many of the metrics used in this evaluation are not especially meaningful on their own, the value often dependent on the magnitude of the observed data, requiring a relative comparison among models forecasting the same targets (1- to 28-day ahead incident hospitalizations for COVID-19) \cite{cramer}.

The COVIDhub-baseline is a naive baseline model created and run by the US COVID-19 Forecast Hub, designed as a simple benchmark for comparison. The baseline makes forecasts using the previous day's incident hospitalization value as the point forecast (and 0.5 quantile for the probabilistic forecast) and using past changes in daily incidence to calibrate uncertainty in its probabilistic forecasts \cite{cramer}. Beating the baseline is typically seen as one marker of a useful model.

The COVIDhub-4\_week\_ensemble is an equally-weighted median ensemble of eligible\footnote{For more information about forecast eligibility, see \cite{ensembles_info}} submitting models for incident hospitalizations at the COVID-19 Forecast Hub, run by the Hub team that includes forecasts up to 28-day (4-week) ahead horizons. This model produces the same 1- to 14-day (1- to 2-week) ahead horizon forecasts as the COVIDhub-ensemble model, with the addition longer horizons \cite{ensembles_info}. These 15- to 28-day (3- to 4-week) horizons were removed from the main ensemble model, as research showed a sharp decline in accuracy in forecasts predicting beyond 2-week aheads \cite{short_horizons}. However, we still include such longer horizons in the analysis for comparison with our models.

Throughout period of analysis, the COVIDhub-4\_week\_ensemble consisted of a shifting list of contributing models due to changing eligibility criteria and submitting models for incident hospitalizations \cite{ensembles_info}. This may have had an impact on shifting accuracy of the ensemble, though the extent is unclear.

The forecasts for the two comparison models are obtained from the Zoltar forecast archive, a research data repository developed by Reich Lab, through the `covidHubUtils` package \cite{zoltar}.


\section{Temporal hierarchical forecasting (THieF)}
Historically, hierarchical forecasting has been utilized in business, economics, operations management, and similar fields. This type of forecasting brings advantages like producing coherent forecasts\footnote{Hierarchical forecasts are coherent if the forecasts at lower levels sum up to those at higher levels. Coherence is not an intrinsic property of all hierarchical forecasts. A lack of coherence causes a dilemma in which it is hard to know which level of forecasts to trust, but a single level must be chosen regardless \cite{hts_talk}.} and increasing prediction accuracy compared to other methods \cite{thief}. Forecast combination has likewise shown favorable outcomes in a myriad of applications. The unique merging of forecast combination with hierarchical forecasting in the THieF methodology has resulted in promising accuracy improvements over more traditional techniques that solely involve hierarchical forecasting \cite{thief}.

Forecast combination that integrates forecasts from different, independent modeling frameworks to create ensembles often results in better infectious disease forecasts compared to those from individual models \cite{ensemble_accuracy}. Likewise, classical statistical models have demonstrated good performance for short-term forecasts in outbreak settings \cite{lauer}. THieF's built-in forecast combination (which shares some key similarities with that of ensemble methods) and underlying ARIMA base-forecaster suggest that this methodology will show accuracy improvements when forecasting COVID-19 incident hospitalizations.

Surveillance data tends to be quite noisy due to reporting artifacts, especially when considering small time scales %\cite{}.
COVID-19 incident hospitalizations suffer a day-of-the-week effect as a result of hospital and clinic operating hours. Fewer hospitalizations are reported on Saturday and Sunday, instead often recorded as part of Monday hospitalizations, since patients are not admitted on weekends. We may better capture true short-term trends by aggregating daily hospitalizations to be weekly or every two weeks, although such smoothing carries some risk of excluding useful information at the daily level. Thus, we are careful to include both daily and weekly aggregation levels in the THieF temporal hierarchy construction for every model. Other higher aggregation levels are included in model schemes to investigate the existence of hidden patterns at higher time scales and if such a pattern could improve model accuracy.

The THieF methodology consists of three steps to produce forecasts from time-series data using a time-based hierarchical structure \cite{thief}. We begin with truth data aggregation to create the temporal hierarchy. The daily confirmed hospitalizations dataset is a time-series consisting of $T = 798$ observations that can be represented as \{$y_t; t = 1, \cdots, 798$\}. These data are then aggregated to different time scales based on the aggregation levels of interest for a particular model (e.g. weekly, every two weeks, every four weeks, etc.\footnote{see THieF models for a full list of time scales considered}). While each aggregation level can be thought of as a separate, univariate time series with different frequencies, together they constitute the temporal hierarchy for one model. Next, we make base forecasts at every level of the hierarchy using an ARIMA model for each level. These forecasts will likely not be coherent, meaning that they predict different trends for the next $h = 28$ days. This disagreement between the forecasts is desirable at this stage, as it may indicate that important patterns have been detected at different time scales. However, through the step of reconciliation (making the forecasts coherent), the information from all of the aggregation levels are combined into forecasts that are expected to be more accurate than any single level. Following this stage, the forecasts from every timescale will agree about the amount of daily incident hospitalizations for the next 1 to 28 days.

\subsection{Data Aggregation and Hierarchy}
Suppose we are interested in constructing a THieF model that combines forecasts from 1-week, 2-week, and 4-week intervals because we believe these time scales each contain important information. Aggregate series for each of these times scales can be constructed from the original daily time series. In fact, any combination of time scales can be chosen for the hierarchy as long as they are factors\footnote{Non-factor levels tend to be more difficult to compute and do not necessarily abide by the following equations.} of the topmost aggregation level (see THieF models for some examples)

$m$ is the number of observations of the smallest timescale (or highest sampling frequency) that aggregate to create a single data point for the highest level. In our example case, $m = 28$ since there are 28 days in 4 weeks.
$k \in$ \{$k_p,\cdot,k_2,k_1$\} can take on multiple values and indicates the particular aggregation level, where $k_p = m$, $k_1 = 1$, and $p$ is the total number of aggregation levels. In our example case, $k \in$ \{$28, 14, 7$\}, where $p = 3$
Athanasopuolos et. al provide a mathematical equation to obtain the data points $y_j^{[k]}$ in each aggregated series, written as \[y_j^{[k]} = \Sigma_{t = t^* + (j-1)k}^{t^*+jk-1} y_t\] for $j = 1, \cdots, \lfloor T/k \rfloor$, $M_k = m/k$ for the aggregated series. We define $t^*= T - \lfloor T/m\rfloor m + 1$ to ensure non-overlapping aggregation
  using our example, we can calculate $t^*= 798 - \lfloor 798/28 \rfloor 28 + 1 = 15$, then we obtain the first data point for the 1-week time series with the equation $y_1^{[7]} = \Sigma_{t = 15 + (0)7}^{15+7-1} y_t = y_{15} + \cdots + y_ {21}$, which is as we expect given that the first $t^*$ observations are removed
However, the index $j$ used above is clearly aggregation-level specific, with no way to translate between levels of the hierarchy. Hence, the authors propose an alternative index $i$ where $i$ is the observation index of the most aggregate series $y_i^{[m]}$ so $j = i$ at that level. A more general form for all levels is $y_j^{[k]} = y_{M_k(i-1)+z}^{[k]}$ for $z = 1, \cdots, M_k$. This new index $i$ is used for the remaining equations and in the figure below.
  Using our example, we can show that the new notation is consistent with the old notation: $y_5^{[7]} = y_{M_7(1)+1}^{[7]} = y_{4+1}^{[7]}$
For each period of the top-level aggregation (2 weeks in our example), we create a column vector of data points


Each aggregation level can be thought of as a univariate time series from which forecasts can be made. Next, a classical statistical model is used to make these base forecasts. These consist of forecasts of up to 84-day ahead incident hospitalizations, depending on the forecasting model (although only up to 28-day ahead forecasts are evaluated).

\subsection{Base Forecasts}
any given date's 1- to 28-day ahead base forecasts represented as \{$\hat{y}_h; h = 1, \cdots, 28$\}
\subsection{Reconciled Forecasts}

\section{Model specifications}
We create three distinct classes of models for comparison during the validation phase: the original THieF models, the THieF ensembles, and the SARIMA models. In total, these three groups consist of twenty-six (26) candidate models.

\subsection{THieF models}
The THieF models are constructed using the THieF methodology outlined in the previous section with different combinations of aggregation levels and data transformations. Seven hierarchical structures and two data transformations result in fourteen (14) original THieF models. Each model is defined (and named) based on the highest level of its temporal hierarchy and the data transformation performed on the forecasts. For example, the simplest of the THieF models is THieF\_1wk-noTransform while the most complicated is THieF\_12wk-4root.

The aggregation levels of any model's temporal hierarchy are daily, plus $f_m$-week periods where $f_m$ is the set of all factors for the highest time scale $m$. Hence, the THieF\_12wk-4root model (as well as the THieF\_12wk-noTransform model) has a temporal hierarchy made up of daily, weekly, 2-weekly, 3-weekly, 4-weekly, 6-weekly, and 12-weekly aggregation levels since $f_{12}=$ \{1, 2, 3, 4, 6, 12\}. Meanwhile, the THieF\_1wk-noTransform model's temporal hierarchy is only made up of daily and weekly aggregation levels ($f_1=$ \{1\}). We chose 1-week, 2-week, 3-week, 4-week, 6-week, 8-week, and 12-week top-level aggregations so as to have a range of time scales to inform the reconciled forecasts.

We explore the seven hierarchy schemes with both no data transformation and a fourth-root data transformation. The latter transform is implemented by first taking the fourth root of the confirmed hospitalizations time series, using the truth data to make forecasts with the THieF methodology, then performing an inverse fourth-power transformation to the resulting forecasts. We selected the fourth-root transformation based on prior research that had shown that had shown accuracy improvements for COVID-19 forecasting due to its variance-stabilizing effect %\cite{}.

\subsection{THieF ensemble models}
The THieF ensemble models\footnote{The equations used to define the ensembles are adapted from Ray, et. al \cite{ensemble_methods}} are seven mean ensembles of the fourteen original THieF models. We formulated six as trained ensembles, in which past performance of the component THieF models is used to assign the components weights, and one as an untrained ensemble, in which all component forecasters are assigned equal weight \cite{ensemble_methods}. Both types of ensembles make quantile forecasts\footnote{Recall that we use the 0.500 quantile forecast as the point forecast.} at level $k$ by combining those of the component forecasts for a particular forecasts date, location, horizon, and quantile combination. (For example, a particular model $m$ might make a 0.500 quantile forecast on December 7, 2020 that predicts the number of 1-day ahead incident hospitalizations in Massachusetts will be 185.) An ensemble quantile forecast can be expressed mathematically as \[q_{l,s,h,k}^{ens} = f(q_{l,s,h,k}^1, \cdots, q_{l,s,h,k}^M),\] where $l$ is the location, $s$ is the forecast date, $h$ is the horizon, and the superscript $m = 1, \cdots, M$ denotes the model \cite{ensemble_methods}.

A trained, weighted mean ensemble calculates its forecast quantiles as \[q_{l,s,h,k}^{ens} = \Sigma_{m=1}^M w_s^m q_{l,s,h,k}^m\], in which the weights are calculated using a sigmoidal transformation of the component models' relative WIS over a rolling window of twelve (12) weeks leading up to the ensemble's forecast date $s$: \[w_s^m=\frac{exp(-\theta \cdot rWIS_s^m)}{\Sigma_{m'=1}^M exp(-\theta) \cdot rWIS_s^{m'}}.\] When the non-negative parameter $\theta = 0$, this becomes an equally weighted ensemble. However, as $\theta$ increases, the weights become more unequal, with the contributions of better-performing component models eclipsing that of worse-performing models. Note that the weights, however, are always non-negative and sum to one \cite{ensemble_methods}.

The THieF ensembles are named with the $\theta$ value used in their construction, except for the untrained ensemble, which is called THieF\_ensemble-mean. We select $\theta$ values \{0, 1, 3, 6.5, 10, 15, 20, 25\} to encompass the range of values taken on by the COVIDhub-trained\_ensemble\footnote{A trained ensemble created by the Forecast Hub that makes short-term forecasts for incident cases, incident and cumulative deaths, and incident hospitalizations. This model makes its forecasts by computing a weighted median of the ten best component forecasts (by WIS) in the 12 weeks leading up to the forecast date. For this model, $\theta$ is allowed to vary by forecast date \cite{ensembles_info, ensemble_methods}.} when predicting incident hospitalization \cite{ensemble_methods}. The one non-integer value (6.5) was selected as the median (and approximate mean) $\theta$ value over the period of analysis for all locations of interest except for the U.S. national level. The corresponding ensemble was named THieF\_ensemble-train6.5.

While the value of $\theta$ generally remains static for each ensemble, all trained ensembles are defined to start with $\theta = 0$ until 12 weeks of forecasts are achieved (12/5/20+12 weeks) that can be used to calculate weights based on prior performance of forecasts.

\subsection{SARIMA models}
The four SARIMA (seasonal ARIMA) models are classic statistical models created using the `auto.arima()` function from the `forecast` package in `R`. This function makes forecasts based on input time series data, and the user may specify seasonality when creating a SARIMA model. Using the original HealthData.Gov confirmed hospitalizations data set, we create two ARIMA models and two SARIMA models, with either no data transformation or a variance-stabilizing fourth root transformation, as described in THieF Models. The SARIMA models both have a seasonality of 7 days (1 week), named sarima\_s7-4root and sarima\_s7-noTransform, while the ARIMA models have no seasonality, named sarima\_s1-4root and sarima\_s1-4root.

The construction of the SARIMA models is intentionally identical to that of the daily and weekly aggregation levels of the original THieF models, though in practice the resulting base forecasts from the THieF models may differ slightly due to the aggregation process to create the temporal hierarchy. We specifically want to use the SARIMA models to help investigate whether it is the THieF methodology's temporal hierarchy and forecast reconciliation that impacts its performance, compared to simply the underlying ARIMA base forecaster. (Recall that research has shown that classic statistical models have shown good performance for outbreak forecasting \cite{lauer}.)

\section{Metrics and Evaluation}
Evaluation of our models is conducted by aggregating the output forecasts in five ways: by forecast state, by location, by pandemic wave, and overall. Despite being made on a daily scale, hospitalization forecasts have been evaluated on a weekly scale by the Forecast Hub, and we follow this convention. For each metric, a single value is obtained by averaging the values across a single week (Tuesday through Monday) \cite{eval_reports}. However, modeling teams have made forecasts on different days (e.g. Saturday vs Sunday vs Monday), leading to a mismatch between what day an $h$-day ahead forecast is predicting for. Hence, for the models we construct, we actually use 3- to 30-day ahead horizons to ensure all forecasts are predicting for the same dates. It should be noted, however, this may result in a slight disadvantage for models whose 3- to 30-day ahead forecasts are being compared to 1- to 28-day ahead forecasts.

We use the same metrics as the Forecast Hub to evaluate model accuracy for both the validation and testing phases. Average mean absolute error (MAE) is used to score point forecasts, while average weighted interval score (WIS), average 50\% prediction interval (PI), and average 95\% PI coverage are used to score probabilistic forecasts. (We refer to these metrics simply as MAE, WIS, and PI coverage going forward for simplicity.)


Relative versions of MAE and WIS are also used to adjust for truth data magnitude to allow for comparison between forecast values of different scales. These relative metrics are calculated as $MAE_{\textrm{model m}} / MAE_{\textrm{baseline}}$ and $WIS_{\textrm{model m}} /WIS_{\textrm{baseline}}$, respectively.

\subsection{Validation phase}
- Serves as time during which to train the candidate models, calibrate the forecasts, and to select the best models to pass on to the testing phase
- All validation phase evaluations are restricted to forecasts with target end dates (date for which the prediction of instant hospitalizations is being made) within the validation phase (July 27, 2020 to October 31, 2021)
- Models passed on: single best stratified by location (states or U.S. national) and model class (regular THieF, SARIMA, THieF ensemble), best regular THieF in last 17 weeks (to account for any models that perform better with more data to train on, especially since they will have a greater amount of data to inform their forecasts during the testing face)

\subsection{Testing phase}
- True time to evaluate the best subset against the two comparison models
- All testing phase evaluations are restricted to forecasts with target end dates (date for which the prediction of instant hospitalizations is being made) within the testing phase (November 1, 2021 to October 2, 2022)

\chapter{Results}
\section{Validation phase}
\begin{table}[h]
\begin{center}
\label{tab:table_validation_us}
% latex table generated in R 4.1.3 by xtable 1.8-4 package
% Tue Mar 21 18:58:31 2023
\begin{tabular}{rlrrrrrr}
  \hline
 & model & wis & mae & cov\_50 & cov\_95 & rwis & rmae \\ 
  \hline
1 & THieF\_6wk-noTransform & 1116.40 & 1756.00 & 0.36 & 0.91 & 0.88 & 0.88 \\ 
  2 & THieF\_ensemble-train3 & 1139.80 & 1808.60 & 0.34 & 0.91 & 0.89 & 0.91 \\ 
  3 & THieF\_ensemble-train1 & 1144.90 & 1814.90 & 0.35 & 0.90 & 0.90 & 0.91 \\ 
  4 & THieF\_ensemble-train6.5 & 1146.60 & 1817.40 & 0.34 & 0.91 & 0.90 & 0.91 \\ 
  5 & THieF\_6wk-4root & 1152.40 & 1765.80 & 0.33 & 0.76 & 0.90 & 0.88 \\ 
  6 & THieF\_ensemble-train10 & 1159.20 & 1834.00 & 0.33 & 0.91 & 0.91 & 0.92 \\ 
  7 & THieF\_ensemble-train15 & 1180.20 & 1866.30 & 0.33 & 0.90 & 0.92 & 0.94 \\ 
  8 & THieF\_ensemble-train20 & 1199.40 & 1897.50 & 0.32 & 0.90 & 0.94 & 0.95 \\ 
  9 & THieF\_ensemble-mean & 1212.70 & 1822.30 & 0.35 & 0.90 & 0.95 & 0.91 \\ 
  10 & THieF\_ensemble-train25 & 1214.30 & 1919.30 & 0.32 & 0.89 & 0.95 & 0.96 \\ 
  11 & THieF\_1wk-noTransform & 1218.50 & 1856.40 & 0.46 & 0.87 & 0.95 & 0.93 \\ 
  12 & THieF\_2wk-4root & 1234.70 & 1941.80 & 0.25 & 0.85 & 0.97 & 0.97 \\ 
  13 & THieF\_1wk-4root & 1252.20 & 1993.10 & 0.30 & 0.85 & 0.98 & 1.00 \\ 
  14 & COVIDhub-baseline & 1276.50 & 1995.90 & 0.64 & 0.99 & 1.00 & 1.00 \\ 
  15 & THieF\_8wk-noTransform & 1277.00 & 1968.90 & 0.38 & 0.85 & 1.00 & 0.99 \\ 
  16 & THieF\_8wk-4root & 1303.60 & 1963.80 & 0.26 & 0.70 & 1.02 & 0.98 \\ 
  17 & THieF\_2wk-noTransform & 1343.10 & 1985.20 & 0.40 & 0.82 & 1.05 & 0.99 \\ 
  18 & THieF\_4wk-4root & 1383.40 & 2058.00 & 0.28 & 0.70 & 1.08 & 1.03 \\ 
  19 & THieF\_3wk-noTransform & 1392.90 & 2008.60 & 0.38 & 0.78 & 1.09 & 1.01 \\ 
  20 & THieF\_4wk-noTransform & 1397.40 & 2038.40 & 0.35 & 0.77 & 1.09 & 1.02 \\ 
  21 & THieF\_3wk-4root & 1401.70 & 2164.80 & 0.26 & 0.80 & 1.10 & 1.08 \\ 
  22 & sarima\_s1-noTransform & 1477.90 & 2162.00 & 0.37 & 0.79 & 1.16 & 1.08 \\ 
  23 & sarima\_s7-noTransform & 1483.20 & 2200.60 & 0.39 & 0.81 & 1.16 & 1.10 \\ 
  24 & sarima\_s1-4root & 1530.20 & 2227.10 & 0.26 & 0.75 & 1.20 & 1.12 \\ 
  25 & THieF\_12wk-noTransform & 1553.50 & 2357.80 & 0.32 & 0.76 & 1.22 & 1.18 \\ 
  26 & sarima\_s7-4root & 1601.50 & 2467.20 & 0.22 & 0.76 & 1.25 & 1.24 \\ 
  27 & THieF\_12wk-4root & 2270.60 & 2039.90 & 0.34 & 0.66 & 1.78 & 1.02 \\ 
   \hline
\end{tabular}\caption{Summary of overall model performance for the US national level, aggregated across forecast week and 1 to 4 week ahead horizons. Prediction interval (PI) coverage rates measure the fraction of times the fraction of times the 50\% and 95\% PIs covered the true values. A well-calibrated model should achieve close to nominal (0.5 and 0.95, respectively) coverage rates. The “rwis” and “rmae” show relative weighted interval score and relative mean absolute error, which compare each model to the baseline. The baseline is defined to have relative metrics of 1. Models that perform better than the baseline have relative WIS or MAE less than 1. The table is ordered by ascending WIS value. Thus two THieF-4wk models show the best performance out of the THieF models, meeting the baseline for both relative WIS and MAE.}
\end{center}
\end{table}

\begin{table}[h]
\begin{center}
\label{tab:table_validation_states}
% latex table generated in R 4.1.3 by xtable 1.8-4 package
% Tue Mar 21 18:58:31 2023
\begin{tabular}{rlrrrrrr}
  \hline
 & model & wis & mae & cov\_50 & cov\_95 & rwis & rmae \\ 
  \hline
1 & THieF\_ensemble-train10 & 30.73 & 47.22 & 0.53 & 0.93 & 0.86 & 0.94 \\ 
  2 & THieF\_ensemble-train15 & 30.74 & 47.28 & 0.52 & 0.93 & 0.86 & 0.95 \\ 
  3 & THieF\_ensemble-train6.5 & 30.77 & 47.20 & 0.53 & 0.94 & 0.86 & 0.94 \\ 
  4 & THieF\_ensemble-train20 & 30.77 & 47.33 & 0.52 & 0.93 & 0.86 & 0.95 \\ 
  5 & THieF\_ensemble-train25 & 30.81 & 47.37 & 0.51 & 0.93 & 0.86 & 0.95 \\ 
  6 & THieF\_ensemble-train3 & 30.90 & 47.25 & 0.54 & 0.94 & 0.86 & 0.95 \\ 
  7 & THieF\_ensemble-train1 & 31.33 & 47.44 & 0.53 & 0.94 & 0.87 & 0.95 \\ 
  8 & THieF\_6wk-4root & 31.65 & 48.11 & 0.46 & 0.91 & 0.88 & 0.96 \\ 
  9 & THieF\_6wk-noTransform & 32.26 & 48.67 & 0.54 & 0.92 & 0.90 & 0.97 \\ 
  10 & THieF\_1wk-4root & 32.68 & 50.06 & 0.48 & 0.91 & 0.91 & 1.00 \\ 
  11 & THieF\_2wk-4root & 32.94 & 50.26 & 0.47 & 0.91 & 0.92 & 1.01 \\ 
  12 & THieF\_2wk-noTransform & 33.26 & 48.03 & 0.57 & 0.90 & 0.93 & 0.96 \\ 
  13 & THieF\_1wk-noTransform & 33.74 & 49.21 & 0.58 & 0.90 & 0.94 & 0.98 \\ 
  14 & THieF\_8wk-4root & 33.84 & 51.14 & 0.43 & 0.89 & 0.94 & 1.02 \\ 
  15 & THieF\_4wk-4root & 34.10 & 51.50 & 0.46 & 0.90 & 0.95 & 1.03 \\ 
  16 & THieF\_ensemble-mean & 34.29 & 47.65 & 0.53 & 0.94 & 0.96 & 0.95 \\ 
  17 & THieF\_3wk-4root & 34.62 & 52.71 & 0.46 & 0.90 & 0.96 & 1.05 \\ 
  18 & THieF\_8wk-noTransform & 34.76 & 52.06 & 0.51 & 0.90 & 0.97 & 1.04 \\ 
  19 & THieF\_3wk-noTransform & 34.96 & 52.08 & 0.54 & 0.90 & 0.97 & 1.04 \\ 
  20 & THieF\_4wk-noTransform & 35.63 & 52.82 & 0.53 & 0.89 & 0.99 & 1.06 \\ 
  21 & COVIDhub-baseline & 35.88 & 49.98 & 0.83 & 0.98 & 1.00 & 1.00 \\ 
  22 & sarima\_s7-noTransform & 36.29 & 51.91 & 0.57 & 0.89 & 1.01 & 1.04 \\ 
  23 & sarima\_s7-4root & 36.35 & 54.55 & 0.46 & 0.90 & 1.01 & 1.09 \\ 
  24 & sarima\_s1-4root & 36.91 & 55.43 & 0.46 & 0.89 & 1.03 & 1.11 \\ 
  25 & THieF\_12wk-noTransform & 39.14 & 59.06 & 0.42 & 0.89 & 1.09 & 1.18 \\ 
  26 & sarima\_s1-noTransform & 39.41 & 56.38 & 0.54 & 0.88 & 1.10 & 1.13 \\ 
  27 & THieF\_12wk-4root & 84.32 & 54.49 & 0.41 & 0.86 & 2.35 & 1.09 \\ 
   \hline
\end{tabular}\caption{Summary of overall model performance averaged over the states level. A similar ordering of models is seen, with the THieF-4wk-4root model performing the best and the THieF-12wk models performing the worst. However, the THieF-8wk models perform better in terms of model ranking for the states level. Additionally, relative WIS sees smaller values and 50\% and 95\% PI coverage see values closer to the nominal levels compared to that at the US national level, although relative MAE is larger with every model performing worse than the baseline.}
\end{center}
\end{table}

\section{Testing phase}

\begin{figure}
\label{fig:figure_forecasts_plot1}
% \includegraphics[width=0.5\textwidth, inner]{overleaf-logo}
\includegraphics{ms_thesis-forecast_models_plot}
\caption{}
\end{figure}

\subsection{Summary of comparison models}
\subsection{Overall model performance}


\begin{table}[h]
\begin{center}
\label{tab:table_overall_us}
% latex table generated in R 4.1.3 by xtable 1.8-4 package
% Tue Mar 21 18:58:48 2023
\begin{tabular}{rlrrrrrr}
  \hline
 & model & wis & mae & cov\_50 & cov\_95 & rwis & rmae \\ 
  \hline
1 & COVIDhub-4\_week\_ensemble & 1401.30 & 2085.10 & 0.63 & 0.91 & 0.75 & 0.83 \\ 
  2 & THieF\_6wk-4root & 1642.50 & 2358.80 & 0.40 & 0.82 & 0.88 & 0.94 \\ 
  3 & THieF\_ensemble-train3 & 1687.40 & 2443.40 & 0.36 & 0.83 & 0.90 & 0.97 \\ 
  4 & THieF\_ensemble-mean & 1700.50 & 2458.90 & 0.39 & 0.84 & 0.91 & 0.98 \\ 
  5 & THieF\_6wk-noTransform & 1727.70 & 2504.80 & 0.46 & 0.86 & 0.92 & 1.00 \\ 
  6 & COVIDhub-baseline & 1867.50 & 2514.10 & 0.71 & 0.85 & 1.00 & 1.00 \\ 
  7 & THieF\_12wk-noTransform & 2074.10 & 2908.30 & 0.37 & 0.86 & 1.11 & 1.16 \\ 
  8 & sarima\_s7-noTransform & 2316.80 & 3042.40 & 0.41 & 0.80 & 1.24 & 1.21 \\ 
   \hline
\end{tabular}\caption{Summary of overall model performance for the US national level, aggregated across forecast week and 1 to 4 week ahead horizons. Prediction interval (PI) coverage rates measure the fraction of times the fraction of times the 50\% and 95\% PIs covered the true values. A well-calibrated model should achieve close to nominal (0.5 and 0.95, respectively) coverage rates. The “rwis” and “rmae” show relative weighted interval score and relative mean absolute error, which compare each model to the baseline. The baseline is defined to have relative metrics of 1. Models that perform better than the baseline have relative WIS or MAE less than 1. The table is ordered by ascending WIS value. Thus two THieF-4wk models show the best performance out of the THieF models, meeting the baseline for both relative WIS and MAE.}
\end{center}
\end{table}

\begin{table}[h]
\begin{center}
\label{tab:table_overall_states}
% latex table generated in R 4.1.3 by xtable 1.8-4 package
% Tue Mar 21 18:58:48 2023
\begin{tabular}{rlrrrrrr}
  \hline
 & model & wis & mae & cov\_50 & cov\_95 & rwis & rmae \\ 
  \hline
1 & COVIDhub-4\_week\_ensemble & 30.18 & 44.87 & 0.64 & 0.94 & 0.72 & 0.79 \\ 
  2 & THieF\_6wk-4root & 34.55 & 50.12 & 0.51 & 0.90 & 0.83 & 0.88 \\ 
  3 & THieF\_ensemble-train3 & 35.62 & 51.57 & 0.56 & 0.91 & 0.86 & 0.91 \\ 
  4 & THieF\_ensemble-mean & 35.91 & 51.80 & 0.57 & 0.91 & 0.86 & 0.91 \\ 
  5 & THieF\_6wk-noTransform & 39.51 & 56.87 & 0.55 & 0.88 & 0.95 & 1.00 \\ 
  6 & COVIDhub-baseline & 41.63 & 56.77 & 0.79 & 0.96 & 1.00 & 1.00 \\ 
  7 & THieF\_12wk-noTransform & 41.66 & 60.97 & 0.51 & 0.88 & 1.00 & 1.07 \\ 
  8 & sarima\_s7-noTransform & 43.54 & 58.94 & 0.60 & 0.87 & 1.05 & 1.04 \\ 
   \hline
\end{tabular}\caption{Summary of overall model performance averaged over the states level. A similar ordering of models is seen, with the THieF-4wk-4root model performing the best and the THieF-12wk models performing the worst. However, the THieF-8wk models perform better in terms of model ranking for the states level. Additionally, relative WIS sees smaller values and 50\% and 95\% PI coverage see values closer to the nominal levels compared to that at the US national level, although relative MAE is larger with every model performing worse than the baseline.}
\end{center}
\end{table}



\subsection{Model performance by horizon week}
\begin{figure}
\label{fig:figure_horizon_plot}
% \includegraphics[width=0.5\textwidth, inner]{overleaf-logo}
\includegraphics{ms_thesis-horizon_plot}
\caption{Average WIS by the horizon week for each model for the US national level and the states level. The figures show similar trends in WIS when forecasts are stratified by horizon week compared to being aggregated across horizon weeks. All models show worse performance relative to the baseline at the 1-week ahead horizon, likely a result of being overconfident at this horizon. Once again, the THieF-4wk-4root model performs the best for WIS. Plots for average MAE are not shown due to their similarity to the average WIS plots.}
\end{figure}

\subsection{Model performance by pandemic wave}
\begin{figure}
\label{fig:figure_wave_plot}
% \includegraphics[width=0.5\textwidth, inner]{overleaf-logo}
\includegraphics{ms_thesis-wave_plot}
\caption{Average WIS by the horizon week and pandemic wave for each model for the US national level and the states level. The figures show similar trends in WIS when forecasts are stratified by horizon week and pandemic wave compared to being stratified only across horizon week. Again, all models show worse performance relative to the baseline at the 1-week ahead horizon likely due to being overconfident at this horizon. Model performance is more similar among the top four models during the Alpha wave compared to during the Winter21 wave. Notably, although the THieF-4wk-4root models largely show the lowest WIS, longer horizons during the Winter21 wave saw better performance by the poor-performing THieF-12wk models. Plots for average MAE are not shown due to their similarity to the average WIS plots.}
\end{figure}


\subsection{Model performance by location}
\begin{figure}
\label{fig:figure_location_plot}
% \includegraphics[width=0.5\textwidth, inner]{overleaf-logo}
\includegraphics{ms_thesis-location_plot}
\caption{Relative WIS by location for each model across all horizons for the entire testing phase. Each square contains the average relative WIS or a model at a particular location, colored based on the relative WIS compare to the baseline. Blue indicates model outperformed the baseline for that location on average, while red indicates the model performed worse than the baseline. Locations are sorted based on cumulative hospitalizations during the testing phase while models are sorted from lowest to highest overall relative WIS, aggregated across states (Table 2). The top three models by relative WIS overall during the testing phase outperformed the baseline in every location.}
\end{figure}

\subsection{Model performance by forecast week}
\begin{figure}
\label{fig:figure_date_plot_us}
% \includegraphics[width=0.5\textwidth, inner]{overleaf-logo}
\includegraphics{ms_thesis-us_date_plot}
\caption{Average h-week ahead WIS and 95\% PI coverage for each model for the US national level. Model performance differs from week to week: the WIS is the lowest during weeks with stable and low incident hospitalization values. The THieF-4wk models show the best performance with lower or comparable WIS to the baseline during the majority of the training period with especially good performance for 1-week ahead horizons. Meanwhile, the THieF-12wk showed lower or comparable WIS to the baseline during high incidence and changed points, which have been historically found difficult to forecast accurately. Plots for h-week ahead MAE are not shown due to their similarity to the WIS plots. All THieF models show very variable coverage rates, especially for a 1 week ahead horizon, only the models without a data transformation showing stable coverage starting around February 2021.}
\end{figure}

\begin{figure}
\label{fig:figure_date_plot_states}
% \includegraphics[width=0.5\textwidth, inner]{overleaf-logo}
\includegraphics{ms_thesis-states_date_plot}
\caption{Average h-week ahead WIS and 95\% PI coverage for each model for the states level. Once again, model performance differs from week to week with the lowest WIS values during weeks with stable and low incident hospitalization values and the THieF-4wk models showing the best performance. However, the THieF-12wk did not show good performance during high incidence and changed points for the states level; instead, THieF-12-4root actually showed some strangely large WIS values during April and May 2021. Plots for h-week ahead MAE are not shown due to their similarity to the WIS plots. One-week ahead coverage rates generally improved as time went on, the best models gaining stable and fairly good coverage rates starting in February 2021. For the 4-week ahead horizon, models without a data transformation showed the best performance and stable coverage starting around February 2021, though THieF-4wk-4root and THieF-8wk-4root also saw decent and fairly stable coverage rates after the downward Winter 21 change point.}
\end{figure}


\chapter{Discussion}
- thief ensembles are best
  - seems like there's no real meaningful difference between mean ensemble and weighted <U+03B8>=3
  - both essentially always beat the baseline for WIS and MAE with all stratified evaluations
  - but don't quite beat the hub-ensemble
- thief\_6wk-4root is next best, then thief\_6wk-noTransform
  - 4root is consistently better and more often beats the baseline for non-coverage metrics
  - noTransform is always as good or better as baseline for non-coverage metrics
- sarima and thief\_12wk-noTransform are the worst
  - tend to be worse than the baseline but except in their specialties:
    - sarima performs very well in static, low-count periods/locations
    - 12wk performs very well in rapidly changing, high-count periods/locations
  - but even if they occasionally beat other models (especially thief ensembles), they tend to be only as good or worse
    - so it seems like ensembling is the way to go to get the best of both worlds

- Limitation of overall metrics being dominated by longer horizons and states with high true hospitalization counts
- Future implementation of log scores (maybe to address the above??)
- Using applications of thief specifically tailored for probabilistic forecasting, not just point forecasting

%% End of body
%%%%%%%%%%%%%%%%%%%%%%%%%%%%%%%%%%%%%%%%%%%%%%%%%%%%%%%%%%%%%%%%%%%%%%%%%%%%%%%

\appendix
\unnumberedchapter{Appendix: Supplemental materials}
%% \chapter{THE FIRST APPENDIX TITLE}
%% ...
%% \chapter{THE SECOND APPENDIX TITLE}
%% ...

\begin{table}[h]
\begin{center}
\label{tab:table_last17_us}
% latex table generated in R 4.1.3 by xtable 1.8-4 package
% Tue Mar 21 18:58:56 2023
\begin{tabular}{rlrrrrrr}
  \hline
 & model & wis & mae & cov\_50 & cov\_95 & rwis & rmae \\ 
  \hline
1 & THieF\_12wk-4root & 1109.40 & 1606.60 & 0.60 & 0.86 & 0.71 & 0.62 \\ 
  2 & THieF\_ensemble-mean & 1147.80 & 1863.50 & 0.37 & 0.91 & 0.74 & 0.71 \\ 
  3 & THieF\_6wk-4root & 1155.80 & 1761.00 & 0.48 & 0.88 & 0.74 & 0.68 \\ 
  4 & THieF\_ensemble-train1 & 1170.20 & 1909.10 & 0.37 & 0.92 & 0.75 & 0.73 \\ 
  5 & THieF\_ensemble-train3 & 1189.00 & 1947.60 & 0.35 & 0.92 & 0.76 & 0.75 \\ 
  6 & THieF\_6wk-noTransform & 1216.00 & 1905.90 & 0.33 & 0.86 & 0.78 & 0.73 \\ 
  7 & THieF\_ensemble-train6.5 & 1229.20 & 2013.30 & 0.34 & 0.94 & 0.79 & 0.77 \\ 
  8 & THieF\_ensemble-train10 & 1274.70 & 2080.70 & 0.34 & 0.94 & 0.82 & 0.80 \\ 
  9 & THieF\_8wk-noTransform & 1293.30 & 2063.20 & 0.35 & 0.86 & 0.83 & 0.79 \\ 
  10 & THieF\_ensemble-train15 & 1341.00 & 2185.10 & 0.33 & 0.94 & 0.86 & 0.84 \\ 
  11 & THieF\_12wk-noTransform & 1386.60 & 1952.40 & 0.55 & 0.77 & 0.89 & 0.75 \\ 
  12 & THieF\_ensemble-train20 & 1399.10 & 2278.40 & 0.33 & 0.92 & 0.90 & 0.87 \\ 
  13 & THieF\_8wk-4root & 1409.10 & 2103.70 & 0.34 & 0.79 & 0.91 & 0.81 \\ 
  14 & THieF\_1wk-noTransform & 1410.30 & 2230.50 & 0.28 & 0.79 & 0.91 & 0.86 \\ 
  15 & THieF\_2wk-4root & 1444.60 & 2287.50 & 0.28 & 0.86 & 0.93 & 0.88 \\ 
  16 & THieF\_ensemble-train25 & 1444.90 & 2348.90 & 0.32 & 0.92 & 0.93 & 0.90 \\ 
  17 & THieF\_4wk-noTransform & 1446.60 & 2130.50 & 0.34 & 0.67 & 0.93 & 0.82 \\ 
  18 & THieF\_3wk-noTransform & 1465.90 & 2126.00 & 0.34 & 0.66 & 0.94 & 0.82 \\ 
  19 & THieF\_2wk-noTransform & 1471.60 & 2181.30 & 0.31 & 0.72 & 0.95 & 0.84 \\ 
  20 & THieF\_3wk-4root & 1490.40 & 2400.00 & 0.36 & 0.90 & 0.96 & 0.92 \\ 
  21 & COVIDhub-baseline & 1555.90 & 2606.50 & 0.37 & 0.97 & 1.00 & 1.00 \\ 
  22 & THieF\_1wk-4root & 1564.10 & 2558.40 & 0.29 & 0.87 & 1.00 & 0.98 \\ 
  23 & THieF\_4wk-4root & 1595.50 & 2428.80 & 0.32 & 0.73 & 1.02 & 0.93 \\ 
  24 & sarima\_s1-noTransform & 1763.40 & 2622.20 & 0.22 & 0.71 & 1.13 & 1.01 \\ 
  25 & sarima\_s1-4root & 1764.90 & 2596.30 & 0.24 & 0.79 & 1.13 & 1.00 \\ 
  26 & sarima\_s7-noTransform & 1958.00 & 2920.00 & 0.15 & 0.63 & 1.26 & 1.12 \\ 
  27 & sarima\_s7-4root & 2198.90 & 3418.20 & 0.18 & 0.75 & 1.41 & 1.31 \\ 
   \hline
\end{tabular}\caption{Summary of overall model performance for the US national level, aggregated across forecast week and 1 to 4 week ahead horizons. Prediction interval (PI) coverage rates measure the fraction of times the fraction of times the 50\% and 95\% PIs covered the true values. A well-calibrated model should achieve close to nominal (0.5 and 0.95, respectively) coverage rates. The “rwis” and “rmae” show relative weighted interval score and relative mean absolute error, which compare each model to the baseline. The baseline is defined to have relative metrics of 1. Models that perform better than the baseline have relative WIS or MAE less than 1. The table is ordered by ascending WIS value. Thus two THieF-4wk models show the best performance out of the THieF models, meeting the baseline for both relative WIS and MAE.}
\end{center}
\end{table}

\begin{table}[h]
\begin{center}
\label{tab:table_last17_states}
% latex table generated in R 4.1.3 by xtable 1.8-4 package
% Tue Mar 21 18:58:56 2023
\begin{tabular}{rlrrrrrr}
  \hline
 & model & wis & mae & cov\_50 & cov\_95 & rwis & rmae \\ 
  \hline
1 & THieF\_ensemble-train6.5 & 32.40 & 49.69 & 0.52 & 0.92 & 0.76 & 0.77 \\ 
  2 & THieF\_ensemble-train10 & 32.41 & 49.81 & 0.52 & 0.93 & 0.76 & 0.77 \\ 
  3 & THieF\_ensemble-train3 & 32.45 & 49.55 & 0.52 & 0.92 & 0.76 & 0.77 \\ 
  4 & THieF\_ensemble-train15 & 32.50 & 50.03 & 0.51 & 0.93 & 0.76 & 0.78 \\ 
  5 & THieF\_ensemble-mean & 32.53 & 49.23 & 0.52 & 0.92 & 0.76 & 0.76 \\ 
  6 & THieF\_ensemble-train1 & 32.54 & 49.49 & 0.52 & 0.92 & 0.76 & 0.77 \\ 
  7 & THieF\_ensemble-train20 & 32.63 & 50.26 & 0.51 & 0.92 & 0.76 & 0.78 \\ 
  8 & THieF\_ensemble-train25 & 32.77 & 50.48 & 0.51 & 0.92 & 0.77 & 0.78 \\ 
  9 & THieF\_6wk-4root & 33.13 & 49.34 & 0.50 & 0.91 & 0.77 & 0.77 \\ 
  10 & THieF\_2wk-4root & 33.73 & 51.33 & 0.50 & 0.91 & 0.79 & 0.80 \\ 
  11 & THieF\_2wk-noTransform & 34.19 & 50.15 & 0.50 & 0.88 & 0.80 & 0.78 \\ 
  12 & THieF\_6wk-noTransform & 34.78 & 51.56 & 0.50 & 0.88 & 0.81 & 0.80 \\ 
  13 & THieF\_1wk-4root & 34.85 & 53.57 & 0.49 & 0.91 & 0.81 & 0.83 \\ 
  14 & THieF\_8wk-4root & 34.91 & 51.66 & 0.48 & 0.90 & 0.82 & 0.80 \\ 
  15 & THieF\_12wk-4root & 34.92 & 50.48 & 0.49 & 0.91 & 0.82 & 0.78 \\ 
  16 & THieF\_8wk-noTransform & 35.45 & 52.57 & 0.49 & 0.88 & 0.83 & 0.82 \\ 
  17 & THieF\_3wk-4root & 36.07 & 54.64 & 0.49 & 0.90 & 0.84 & 0.85 \\ 
  18 & THieF\_4wk-4root & 36.66 & 54.82 & 0.48 & 0.89 & 0.86 & 0.85 \\ 
  19 & THieF\_3wk-noTransform & 36.98 & 54.12 & 0.48 & 0.87 & 0.86 & 0.84 \\ 
  20 & THieF\_1wk-noTransform & 37.21 & 54.92 & 0.47 & 0.86 & 0.87 & 0.85 \\ 
  21 & THieF\_4wk-noTransform & 38.33 & 56.40 & 0.47 & 0.86 & 0.90 & 0.88 \\ 
  22 & THieF\_12wk-noTransform & 38.49 & 53.65 & 0.49 & 0.85 & 0.90 & 0.83 \\ 
  23 & sarima\_s1-4root & 41.55 & 62.71 & 0.44 & 0.88 & 0.97 & 0.97 \\ 
  24 & COVIDhub-baseline & 42.78 & 64.41 & 0.72 & 0.97 & 1.00 & 1.00 \\ 
  25 & sarima\_s7-4root & 43.72 & 64.94 & 0.45 & 0.88 & 1.02 & 1.01 \\ 
  26 & sarima\_s7-noTransform & 43.81 & 62.42 & 0.45 & 0.84 & 1.02 & 0.97 \\ 
  27 & sarima\_s1-noTransform & 46.51 & 67.02 & 0.42 & 0.83 & 1.09 & 1.04 \\ 
   \hline
\end{tabular}\caption{Summary of overall model performance averaged over the states level. A similar ordering of models is seen, with the THieF-4wk-4root model performing the best and the THieF-12wk models performing the worst. However, the THieF-8wk models perform better in terms of model ranking for the states level. Additionally, relative WIS sees smaller values and 50\% and 95\% PI coverage see values closer to the nominal levels compared to that at the US national level, although relative MAE is larger with every model performing worse than the baseline.}
\end{center}
\end{table}


\begin{table}[h]
\begin{center}
\label{tab:table_horizon_us}
% latex table generated in R 4.1.3 by xtable 1.8-4 package
% Tue Mar 21 18:58:56 2023
\begin{tabular}{rlrrrrrrr}
  \hline
 & model & horizon\_wk & wis & mae & cov\_50 & cov\_95 & rwis & rmae \\ 
  \hline
1 & COVIDhub-4\_week\_ensemble & 1.00 & 659.80 & 966.80 & 0.68 & 0.94 & 0.69 & 0.77 \\ 
  2 & THieF\_ensemble-train3 & 1.00 & 909.60 & 1348.80 & 0.37 & 0.86 & 0.95 & 1.07 \\ 
  3 & THieF\_ensemble-mean & 1.00 & 930.90 & 1372.90 & 0.38 & 0.86 & 0.98 & 1.09 \\ 
  4 & COVIDhub-baseline & 1.00 & 952.20 & 1258.80 & 0.73 & 0.92 & 1.00 & 1.00 \\ 
  5 & THieF\_6wk-4root & 1.00 & 1007.40 & 1419.10 & 0.45 & 0.79 & 1.06 & 1.13 \\ 
  6 & THieF\_6wk-noTransform & 1.00 & 1028.70 & 1467.40 & 0.46 & 0.88 & 1.08 & 1.17 \\ 
  7 & sarima\_s7-noTransform & 1.00 & 1143.40 & 1533.90 & 0.32 & 0.81 & 1.20 & 1.22 \\ 
  8 & THieF\_12wk-noTransform & 1.00 & 1272.40 & 1793.60 & 0.40 & 0.88 & 1.34 & 1.42 \\ 
  9 & COVIDhub-4\_week\_ensemble & 2.00 & 1156.20 & 1694.70 & 0.66 & 0.92 & 0.72 & 0.80 \\ 
  10 & THieF\_6wk-4root & 2.00 & 1413.20 & 2029.00 & 0.42 & 0.85 & 0.88 & 0.96 \\ 
  11 & THieF\_ensemble-train3 & 2.00 & 1417.10 & 2046.90 & 0.39 & 0.84 & 0.88 & 0.97 \\ 
  12 & THieF\_ensemble-mean & 2.00 & 1436.90 & 2070.80 & 0.40 & 0.85 & 0.90 & 0.98 \\ 
  13 & THieF\_6wk-noTransform & 2.00 & 1528.50 & 2177.00 & 0.47 & 0.86 & 0.95 & 1.03 \\ 
  14 & COVIDhub-baseline & 2.00 & 1604.90 & 2113.30 & 0.74 & 0.85 & 1.00 & 1.00 \\ 
  15 & THieF\_12wk-noTransform & 2.00 & 1858.50 & 2558.30 & 0.38 & 0.85 & 1.16 & 1.21 \\ 
  16 & sarima\_s7-noTransform & 2.00 & 1923.10 & 2536.30 & 0.41 & 0.82 & 1.20 & 1.20 \\ 
  17 & COVIDhub-4\_week\_ensemble & 3.00 & 1670.60 & 2528.80 & 0.60 & 0.90 & 0.75 & 0.84 \\ 
  18 & THieF\_6wk-4root & 3.00 & 1867.40 & 2715.30 & 0.37 & 0.84 & 0.84 & 0.91 \\ 
  19 & THieF\_ensemble-train3 & 3.00 & 1979.60 & 2845.10 & 0.35 & 0.81 & 0.89 & 0.95 \\ 
  20 & THieF\_6wk-noTransform & 3.00 & 1989.00 & 2883.60 & 0.46 & 0.85 & 0.90 & 0.96 \\ 
  21 & THieF\_ensemble-mean & 3.00 & 1989.00 & 2860.10 & 0.39 & 0.83 & 0.90 & 0.95 \\ 
  22 & COVIDhub-baseline & 3.00 & 2217.00 & 2993.70 & 0.71 & 0.83 & 1.00 & 1.00 \\ 
  23 & THieF\_12wk-noTransform & 3.00 & 2369.90 & 3313.80 & 0.37 & 0.86 & 1.07 & 1.11 \\ 
  24 & sarima\_s7-noTransform & 3.00 & 2738.40 & 3581.40 & 0.46 & 0.80 & 1.24 & 1.20 \\ 
  25 & COVIDhub-4\_week\_ensemble & 4.00 & 2173.40 & 3232.70 & 0.56 & 0.89 & 0.79 & 0.85 \\ 
  26 & THieF\_6wk-4root & 4.00 & 2330.00 & 3341.70 & 0.36 & 0.80 & 0.84 & 0.88 \\ 
  27 & THieF\_6wk-noTransform & 4.00 & 2414.90 & 3567.40 & 0.43 & 0.85 & 0.88 & 0.94 \\ 
  28 & THieF\_ensemble-train3 & 4.00 & 2501.10 & 3615.30 & 0.32 & 0.80 & 0.91 & 0.96 \\ 
  29 & THieF\_ensemble-mean & 4.00 & 2502.20 & 3613.60 & 0.38 & 0.80 & 0.91 & 0.95 \\ 
  30 & COVIDhub-baseline & 4.00 & 2761.40 & 3782.10 & 0.67 & 0.80 & 1.00 & 1.00 \\ 
  31 & THieF\_12wk-noTransform & 4.00 & 2852.50 & 4049.10 & 0.31 & 0.86 & 1.03 & 1.07 \\ 
  32 & sarima\_s7-noTransform & 4.00 & 3549.40 & 4630.40 & 0.45 & 0.76 & 1.28 & 1.22 \\ 
   \hline
\end{tabular}\caption{Summary of model performance stratified by horizon week for the US national level, aggregated across forecast week. The results from this table were used to create Figure 4 that summarized average WIS by horizon week for the US national level and the states level. The table is ordered by ascending horizon week and WIS, with THieF-4wk-4root having the lowest WIS except for the 4 week ahead horizon. The baseline consistently does not have the best WIS.}
\end{center}
\end{table}

\begin{table}[h]
\begin{center}
\label{tab:table_horizon_states}
% latex table generated in R 4.1.3 by xtable 1.8-4 package
% Tue Mar 21 18:58:56 2023
\begin{tabular}{rlrrrrrrr}
  \hline
 & model & horizon\_wk & wis & mae & cov\_50 & cov\_95 & rwis & rmae \\ 
  \hline
1 & COVIDhub-4\_week\_ensemble & 1.00 & 16.70 & 25.11 & 0.65 & 0.96 & 0.70 & 0.79 \\ 
  2 & THieF\_ensemble-train3 & 1.00 & 20.32 & 30.04 & 0.55 & 0.93 & 0.85 & 0.94 \\ 
  3 & THieF\_ensemble-mean & 1.00 & 20.76 & 30.52 & 0.56 & 0.93 & 0.87 & 0.96 \\ 
  4 & THieF\_6wk-4root & 1.00 & 22.50 & 32.97 & 0.49 & 0.89 & 0.94 & 1.03 \\ 
  5 & sarima\_s7-noTransform & 1.00 & 23.11 & 31.92 & 0.57 & 0.89 & 0.97 & 1.00 \\ 
  6 & COVIDhub-baseline & 1.00 & 23.92 & 31.87 & 0.80 & 0.97 & 1.00 & 1.00 \\ 
  7 & THieF\_6wk-noTransform & 1.00 & 25.17 & 36.41 & 0.52 & 0.89 & 1.05 & 1.14 \\ 
  8 & THieF\_12wk-noTransform & 1.00 & 29.17 & 42.20 & 0.49 & 0.89 & 1.22 & 1.32 \\ 
  9 & COVIDhub-4\_week\_ensemble & 2.00 & 25.42 & 37.90 & 0.66 & 0.95 & 0.69 & 0.77 \\ 
  10 & THieF\_6wk-4root & 2.00 & 30.10 & 43.94 & 0.52 & 0.91 & 0.81 & 0.89 \\ 
  11 & THieF\_ensemble-train3 & 2.00 & 30.66 & 44.49 & 0.56 & 0.91 & 0.83 & 0.90 \\ 
  12 & THieF\_ensemble-mean & 2.00 & 31.04 & 44.80 & 0.58 & 0.91 & 0.84 & 0.90 \\ 
  13 & THieF\_6wk-noTransform & 2.00 & 34.99 & 50.27 & 0.56 & 0.88 & 0.95 & 1.02 \\ 
  14 & COVIDhub-baseline & 2.00 & 36.95 & 49.48 & 0.80 & 0.96 & 1.00 & 1.00 \\ 
  15 & sarima\_s7-noTransform & 2.00 & 37.32 & 50.28 & 0.60 & 0.87 & 1.01 & 1.02 \\ 
  16 & THieF\_12wk-noTransform & 2.00 & 37.35 & 54.31 & 0.53 & 0.88 & 1.01 & 1.10 \\ 
  17 & COVIDhub-4\_week\_ensemble & 3.00 & 35.22 & 52.52 & 0.64 & 0.94 & 0.73 & 0.79 \\ 
  18 & THieF\_6wk-4root & 3.00 & 38.98 & 56.36 & 0.52 & 0.90 & 0.80 & 0.85 \\ 
  19 & THieF\_ensemble-train3 & 3.00 & 41.60 & 59.88 & 0.56 & 0.90 & 0.86 & 0.90 \\ 
  20 & THieF\_ensemble-mean & 3.00 & 41.85 & 60.07 & 0.58 & 0.90 & 0.86 & 0.91 \\ 
  21 & THieF\_6wk-noTransform & 3.00 & 45.11 & 64.83 & 0.56 & 0.87 & 0.93 & 0.98 \\ 
  22 & THieF\_12wk-noTransform & 3.00 & 46.50 & 68.16 & 0.52 & 0.87 & 0.96 & 1.03 \\ 
  23 & COVIDhub-baseline & 3.00 & 48.54 & 66.30 & 0.79 & 0.95 & 1.00 & 1.00 \\ 
  24 & sarima\_s7-noTransform & 3.00 & 51.31 & 69.07 & 0.60 & 0.86 & 1.06 & 1.04 \\ 
  25 & COVIDhub-4\_week\_ensemble & 4.00 & 44.39 & 65.43 & 0.63 & 0.93 & 0.76 & 0.80 \\ 
  26 & THieF\_6wk-4root & 4.00 & 47.53 & 68.48 & 0.52 & 0.88 & 0.81 & 0.84 \\ 
  27 & THieF\_ensemble-train3 & 4.00 & 51.04 & 73.47 & 0.55 & 0.89 & 0.88 & 0.90 \\ 
  28 & THieF\_ensemble-mean & 4.00 & 51.11 & 73.39 & 0.56 & 0.88 & 0.88 & 0.90 \\ 
  29 & THieF\_6wk-noTransform & 4.00 & 53.80 & 77.46 & 0.55 & 0.86 & 0.92 & 0.95 \\ 
  30 & THieF\_12wk-noTransform & 4.00 & 54.52 & 80.61 & 0.50 & 0.87 & 0.93 & 0.99 \\ 
  31 & COVIDhub-baseline & 4.00 & 58.35 & 81.23 & 0.77 & 0.95 & 1.00 & 1.00 \\ 
  32 & sarima\_s7-noTransform & 4.00 & 63.89 & 86.44 & 0.60 & 0.85 & 1.09 & 1.06 \\ 
   \hline
\end{tabular}\caption{Summary of model performance stratified by horizon week for the states level, aggregated across forecast week. The results from this table were used to create Figure 4 that summarized average WIS by horizon week for the US national level and the states level. The table is ordered by ascending horizon week and WIS, with THieF-4wk-4root having the lowest WIS except for the 3 week (although it is second lowest) and 4 week ahead horizon. The baseline again consistently does not have the best WIS, performing progressively worse compared to the other models for each additional week ahead horizon.}
\end{center}
\end{table}



\begin{table}[h]
\begin{center}
\label{tab:table_wave_us}
% latex table generated in R 4.1.3 by xtable 1.8-4 package
% Tue Mar 21 18:58:56 2023
\begin{longtable}{rlrlrrrrrr}
  \hline
 & model & horizon\_wk & wave & wis & mae & cov\_50 & cov\_95 & rwis & rmae \\ 
  \hline
1 & COVIDhub-4\_week\_ensemble & 1.00 & ba4\_ba5 & 277.80 & 378.70 & 0.89 & 1.00 & 0.62 & 0.78 \\ 
  2 & THieF\_6wk-4root & 1.00 & ba4\_ba5 & 350.40 & 542.90 & 0.60 & 0.99 & 0.78 & 1.11 \\ 
  3 & THieF\_ensemble-train3 & 1.00 & ba4\_ba5 & 399.10 & 685.30 & 0.46 & 1.00 & 0.88 & 1.41 \\ 
  4 & THieF\_ensemble-mean & 1.00 & ba4\_ba5 & 433.60 & 740.30 & 0.49 & 1.00 & 0.96 & 1.52 \\ 
  5 & COVIDhub-baseline & 1.00 & ba4\_ba5 & 450.80 & 486.90 & 0.92 & 1.00 & 1.00 & 1.00 \\ 
  6 & THieF\_6wk-noTransform & 1.00 & ba4\_ba5 & 516.50 & 822.50 & 0.60 & 1.00 & 1.15 & 1.69 \\ 
  7 & sarima\_s7-noTransform & 1.00 & ba4\_ba5 & 517.80 & 828.60 & 0.41 & 0.99 & 1.15 & 1.70 \\ 
  8 & THieF\_12wk-noTransform & 1.00 & ba4\_ba5 & 790.70 & 1319.20 & 0.35 & 1.00 & 1.75 & 2.71 \\ 
  9 & COVIDhub-4\_week\_ensemble & 2.00 & ba4\_ba5 & 411.40 & 572.50 & 0.85 & 1.00 & 0.63 & 0.77 \\ 
  10 & THieF\_6wk-4root & 2.00 & ba4\_ba5 & 518.80 & 813.00 & 0.61 & 1.00 & 0.79 & 1.10 \\ 
  11 & THieF\_ensemble-train3 & 2.00 & ba4\_ba5 & 578.60 & 966.60 & 0.49 & 1.00 & 0.88 & 1.30 \\ 
  12 & THieF\_ensemble-mean & 2.00 & ba4\_ba5 & 621.80 & 1030.30 & 0.50 & 1.00 & 0.95 & 1.39 \\ 
  13 & COVIDhub-baseline & 2.00 & ba4\_ba5 & 654.20 & 741.00 & 0.98 & 1.00 & 1.00 & 1.00 \\ 
  14 & THieF\_6wk-noTransform & 2.00 & ba4\_ba5 & 727.80 & 1145.90 & 0.59 & 1.00 & 1.11 & 1.55 \\ 
  15 & sarima\_s7-noTransform & 2.00 & ba4\_ba5 & 729.70 & 1169.60 & 0.53 & 1.00 & 1.11 & 1.58 \\ 
  16 & THieF\_12wk-noTransform & 2.00 & ba4\_ba5 & 1012.60 & 1677.50 & 0.38 & 1.00 & 1.55 & 2.26 \\ 
  17 & COVIDhub-4\_week\_ensemble & 3.00 & ba4\_ba5 & 563.30 & 824.70 & 0.77 & 1.00 & 0.68 & 0.79 \\ 
  18 & THieF\_6wk-4root & 3.00 & ba4\_ba5 & 683.90 & 1069.10 & 0.59 & 1.00 & 0.83 & 1.03 \\ 
  19 & THieF\_ensemble-train3 & 3.00 & ba4\_ba5 & 784.00 & 1311.50 & 0.45 & 1.00 & 0.95 & 1.26 \\ 
  20 & COVIDhub-baseline & 3.00 & ba4\_ba5 & 825.00 & 1038.90 & 1.00 & 1.00 & 1.00 & 1.00 \\ 
  21 & THieF\_ensemble-mean & 3.00 & ba4\_ba5 & 825.00 & 1371.30 & 0.50 & 1.00 & 1.00 & 1.32 \\ 
  22 & THieF\_6wk-noTransform & 3.00 & ba4\_ba5 & 911.00 & 1464.40 & 0.59 & 1.00 & 1.10 & 1.41 \\ 
  23 & sarima\_s7-noTransform & 3.00 & ba4\_ba5 & 953.50 & 1514.20 & 0.57 & 1.00 & 1.16 & 1.46 \\ 
  24 & THieF\_12wk-noTransform & 3.00 & ba4\_ba5 & 1248.70 & 2097.10 & 0.38 & 1.00 & 1.51 & 2.02 \\ 
  25 & COVIDhub-4\_week\_ensemble & 4.00 & ba4\_ba5 & 698.90 & 1034.30 & 0.77 & 1.00 & 0.72 & 0.78 \\ 
  26 & THieF\_6wk-4root & 4.00 & ba4\_ba5 & 801.90 & 1210.30 & 0.64 & 1.00 & 0.82 & 0.92 \\ 
  27 & THieF\_ensemble-train3 & 4.00 & ba4\_ba5 & 969.80 & 1585.10 & 0.48 & 1.00 & 1.00 & 1.20 \\ 
  28 & COVIDhub-baseline & 4.00 & ba4\_ba5 & 972.00 & 1320.50 & 1.00 & 1.00 & 1.00 & 1.00 \\ 
  29 & THieF\_ensemble-mean & 4.00 & ba4\_ba5 & 1010.80 & 1651.70 & 0.53 & 1.00 & 1.04 & 1.25 \\ 
  30 & THieF\_6wk-noTransform & 4.00 & ba4\_ba5 & 1034.00 & 1631.70 & 0.61 & 1.00 & 1.06 & 1.24 \\ 
  31 & sarima\_s7-noTransform & 4.00 & ba4\_ba5 & 1169.60 & 1821.80 & 0.60 & 1.00 & 1.20 & 1.38 \\ 
  32 & THieF\_12wk-noTransform & 4.00 & ba4\_ba5 & 1501.80 & 2581.80 & 0.34 & 1.00 & 1.54 & 1.96 \\ 
  33 & THieF\_ensemble-mean & 1.00 & delta & 239.10 & 305.40 & 0.86 & 1.00 & 0.16 & 0.15 \\ 
  34 & THieF\_ensemble-train3 & 1.00 & delta & 247.90 & 326.80 & 0.86 & 1.00 & 0.17 & 0.16 \\ 
  35 & sarima\_s7-noTransform & 1.00 & delta & 270.90 & 454.30 & 0.29 & 1.00 & 0.18 & 0.22 \\ 
  36 & THieF\_6wk-4root & 1.00 & delta & 327.60 & 536.40 & 0.71 & 1.00 & 0.22 & 0.26 \\ 
  37 & THieF\_6wk-noTransform & 1.00 & delta & 332.50 & 563.70 & 0.43 & 1.00 & 0.22 & 0.27 \\ 
  38 & THieF\_12wk-noTransform & 1.00 & delta & 381.80 & 523.00 & 1.00 & 1.00 & 0.26 & 0.25 \\ 
  39 & THieF\_6wk-noTransform & 2.00 & delta & 286.50 & 260.10 & 1.00 & 1.00 & 0.11 & 0.07 \\ 
  40 & sarima\_s7-noTransform & 2.00 & delta & 290.70 & 240.00 & 1.00 & 1.00 & 0.11 & 0.07 \\ 
  41 & THieF\_ensemble-mean & 2.00 & delta & 313.60 & 215.30 & 1.00 & 1.00 & 0.12 & 0.06 \\ 
  42 & THieF\_ensemble-train3 & 2.00 & delta & 316.20 & 207.60 & 1.00 & 1.00 & 0.12 & 0.06 \\ 
  43 & THieF\_6wk-4root & 2.00 & delta & 324.40 & 280.40 & 1.00 & 1.00 & 0.13 & 0.08 \\ 
  44 & THieF\_12wk-noTransform & 2.00 & delta & 487.90 & 404.70 & 1.00 & 1.00 & 0.19 & 0.12 \\ 
  45 & THieF\_6wk-noTransform & 3.00 & delta & 396.00 & 511.10 & 0.86 & 1.00 & 0.11 & 0.10 \\ 
  46 & sarima\_s7-noTransform & 3.00 & delta & 486.30 & 701.80 & 1.00 & 1.00 & 0.14 & 0.14 \\ 
  47 & THieF\_ensemble-mean & 3.00 & delta & 509.60 & 800.20 & 1.00 & 1.00 & 0.14 & 0.16 \\ 
  48 & THieF\_ensemble-train3 & 3.00 & delta & 509.60 & 797.20 & 1.00 & 1.00 & 0.14 & 0.16 \\ 
  49 & THieF\_6wk-4root & 3.00 & delta & 547.90 & 979.90 & 0.86 & 1.00 & 0.15 & 0.20 \\ 
  50 & THieF\_12wk-noTransform & 3.00 & delta & 585.80 & 276.30 & 1.00 & 1.00 & 0.16 & 0.06 \\ 
  51 & THieF\_6wk-noTransform & 4.00 & delta & 530.60 & 831.20 & 0.71 & 1.00 & 0.12 & 0.14 \\ 
  52 & THieF\_ensemble-mean & 4.00 & delta & 691.00 & 1128.00 & 0.86 & 1.00 & 0.16 & 0.18 \\ 
  53 & THieF\_ensemble-train3 & 4.00 & delta & 691.80 & 1135.10 & 0.86 & 1.00 & 0.16 & 0.19 \\ 
  54 & THieF\_12wk-noTransform & 4.00 & delta & 710.80 & 555.60 & 1.00 & 1.00 & 0.16 & 0.09 \\ 
  55 & sarima\_s7-noTransform & 4.00 & delta & 745.20 & 1232.60 & 0.71 & 1.00 & 0.17 & 0.20 \\ 
  56 & THieF\_6wk-4root & 4.00 & delta & 801.60 & 1496.10 & 0.43 & 1.00 & 0.18 & 0.25 \\ 
  57 & COVIDhub-4\_week\_ensemble & 1.00 & omicron & 1067.90 & 1595.10 & 0.46 & 0.87 & 0.72 & 0.77 \\ 
  58 & COVIDhub-baseline & 1.00 & omicron & 1487.90 & 2083.40 & 0.53 & 0.83 & 1.00 & 1.00 \\ 
  59 & THieF\_ensemble-train3 & 1.00 & omicron & 1509.90 & 2136.20 & 0.24 & 0.69 & 3.35 & 4.39 \\ 
  60 & THieF\_ensemble-mean & 1.00 & omicron & 1517.90 & 2127.80 & 0.24 & 0.69 & 3.37 & 4.37 \\ 
  61 & THieF\_6wk-noTransform & 1.00 & omicron & 1632.50 & 2228.70 & 0.32 & 0.73 & 3.62 & 4.58 \\ 
  62 & THieF\_6wk-4root & 1.00 & omicron & 1772.10 & 2437.80 & 0.26 & 0.54 & 3.93 & 5.01 \\ 
  63 & THieF\_12wk-noTransform & 1.00 & omicron & 1850.90 & 2381.20 & 0.43 & 0.75 & 4.11 & 4.89 \\ 
  64 & sarima\_s7-noTransform & 1.00 & omicron & 1881.90 & 2370.80 & 0.22 & 0.60 & 4.17 & 4.87 \\ 
  65 & COVIDhub-4\_week\_ensemble & 2.00 & omicron & 1919.40 & 2844.90 & 0.47 & 0.84 & 0.74 & 0.81 \\ 
  66 & THieF\_ensemble-mean & 2.00 & omicron & 2361.20 & 3270.00 & 0.27 & 0.68 & 3.61 & 4.41 \\ 
  67 & THieF\_ensemble-train3 & 2.00 & omicron & 2365.40 & 3287.90 & 0.26 & 0.67 & 3.62 & 4.44 \\ 
  68 & THieF\_6wk-4root & 2.00 & omicron & 2421.10 & 3411.40 & 0.19 & 0.68 & 3.70 & 4.60 \\ 
  69 & THieF\_6wk-noTransform & 2.00 & omicron & 2442.90 & 3368.80 & 0.31 & 0.71 & 3.73 & 4.55 \\ 
  70 & COVIDhub-baseline & 2.00 & omicron & 2579.10 & 3519.70 & 0.50 & 0.70 & 1.00 & 1.00 \\ 
  71 & THieF\_12wk-noTransform & 2.00 & omicron & 2827.30 & 3600.00 & 0.36 & 0.69 & 4.32 & 4.86 \\ 
  72 & sarima\_s7-noTransform & 2.00 & omicron & 3275.90 & 4104.90 & 0.25 & 0.63 & 5.01 & 5.54 \\ 
  73 & COVIDhub-4\_week\_ensemble & 3.00 & omicron & 2757.20 & 4201.20 & 0.43 & 0.80 & 0.77 & 0.85 \\ 
  74 & THieF\_6wk-4root & 3.00 & omicron & 3141.50 & 4483.10 & 0.12 & 0.67 & 3.81 & 4.32 \\ 
  75 & THieF\_6wk-noTransform & 3.00 & omicron & 3167.40 & 4447.50 & 0.30 & 0.69 & 3.84 & 4.28 \\ 
  76 & THieF\_ensemble-mean & 3.00 & omicron & 3250.50 & 4481.10 & 0.25 & 0.66 & 3.94 & 4.31 \\ 
  77 & THieF\_ensemble-train3 & 3.00 & omicron & 3273.10 & 4511.80 & 0.22 & 0.62 & 3.97 & 4.34 \\ 
  78 & COVIDhub-baseline & 3.00 & omicron & 3583.00 & 4912.10 & 0.43 & 0.67 & 1.00 & 1.00 \\ 
  79 & THieF\_12wk-noTransform & 3.00 & omicron & 3601.40 & 4700.20 & 0.34 & 0.72 & 4.37 & 4.52 \\ 
  80 & sarima\_s7-noTransform & 3.00 & omicron & 4672.10 & 5833.20 & 0.32 & 0.58 & 5.66 & 5.62 \\ 
  81 & COVIDhub-4\_week\_ensemble & 4.00 & omicron & 3556.30 & 5294.60 & 0.37 & 0.79 & 0.80 & 0.87 \\ 
  82 & THieF\_6wk-noTransform & 4.00 & omicron & 3854.40 & 5589.80 & 0.25 & 0.70 & 3.96 & 4.23 \\ 
  83 & THieF\_6wk-4root & 4.00 & omicron & 3897.80 & 5515.50 & 0.08 & 0.58 & 4.01 & 4.18 \\ 
  84 & THieF\_ensemble-mean & 4.00 & omicron & 4046.90 & 5650.10 & 0.20 & 0.60 & 4.16 & 4.28 \\ 
  85 & THieF\_ensemble-train3 & 4.00 & omicron & 4084.80 & 5718.60 & 0.14 & 0.58 & 4.20 & 4.33 \\ 
  86 & THieF\_12wk-noTransform & 4.00 & omicron & 4274.30 & 5646.60 & 0.25 & 0.71 & 4.40 & 4.28 \\ 
  87 & COVIDhub-baseline & 4.00 & omicron & 4439.60 & 6090.80 & 0.36 & 0.61 & 1.00 & 1.00 \\ 
  88 & sarima\_s7-noTransform & 4.00 & omicron & 6010.40 & 7538.60 & 0.29 & 0.51 & 6.18 & 5.71 \\ 
   \hline
\hline
\end{longtable}\caption{Summary of model performance stratified by horizon week and pandemic waves for the US national level, aggregated across forecast week. The results from this table were used to create Figure 5 that summarized average WIS by horizon week and pandemic wave for the US national level and the states level. The table is ordered by ascending pandemic wave, horizon week, and WIS, with THieF-4wk-4root having the lowest WIS for all horizons of the alpha wave and shorter horizons of the Winter 21 wave. The THieF-12wk-noTrans and THieF-12wk-4root models have the lowest and second lowest WIS for 3 week ahead and 4 week ahead horizons during the Winter 21 wave. The baseline consistently does not have the best WIS.}
\end{center}
\end{table}

\begin{table}[h]
\begin{center}
\label{tab:table_wave_states}
% latex table generated in R 4.1.3 by xtable 1.8-4 package
% Tue Mar 21 18:58:56 2023
\begin{longtable}{rlrlrrrrrr}
  \hline
 & model & horizon\_wk & wave & wis & mae & cov\_50 & cov\_95 & rwis & rmae \\ 
  \hline
1 & COVIDhub-4\_week\_ensemble & 1.00 & ba4\_ba5 & 8.91 & 13.43 & 0.72 & 0.99 & 0.61 & 0.81 \\ 
  2 & THieF\_6wk-4root & 1.00 & ba4\_ba5 & 9.95 & 14.83 & 0.62 & 0.97 & 0.69 & 0.89 \\ 
  3 & THieF\_ensemble-train3 & 1.00 & ba4\_ba5 & 10.48 & 15.62 & 0.67 & 0.99 & 0.72 & 0.94 \\ 
  4 & THieF\_ensemble-mean & 1.00 & ba4\_ba5 & 11.10 & 16.47 & 0.69 & 0.99 & 0.77 & 0.99 \\ 
  5 & sarima\_s7-noTransform & 1.00 & ba4\_ba5 & 12.81 & 19.00 & 0.68 & 0.98 & 0.88 & 1.14 \\ 
  6 & THieF\_6wk-noTransform & 1.00 & ba4\_ba5 & 13.93 & 20.54 & 0.70 & 0.99 & 0.96 & 1.24 \\ 
  7 & COVIDhub-baseline & 1.00 & ba4\_ba5 & 14.49 & 16.59 & 0.92 & 1.00 & 1.00 & 1.00 \\ 
  8 & THieF\_12wk-noTransform & 1.00 & ba4\_ba5 & 18.37 & 28.32 & 0.61 & 0.99 & 1.27 & 1.71 \\ 
  9 & COVIDhub-4\_week\_ensemble & 2.00 & ba4\_ba5 & 11.85 & 17.64 & 0.76 & 1.00 & 0.60 & 0.81 \\ 
  10 & THieF\_6wk-4root & 2.00 & ba4\_ba5 & 13.28 & 19.07 & 0.69 & 0.99 & 0.68 & 0.88 \\ 
  11 & THieF\_ensemble-train3 & 2.00 & ba4\_ba5 & 14.15 & 20.34 & 0.74 & 0.99 & 0.72 & 0.94 \\ 
  12 & THieF\_ensemble-mean & 2.00 & ba4\_ba5 & 14.86 & 21.20 & 0.75 & 0.99 & 0.76 & 0.98 \\ 
  13 & sarima\_s7-noTransform & 2.00 & ba4\_ba5 & 17.46 & 25.07 & 0.75 & 0.99 & 0.89 & 1.16 \\ 
  14 & THieF\_6wk-noTransform & 2.00 & ba4\_ba5 & 17.92 & 25.69 & 0.77 & 1.00 & 0.91 & 1.19 \\ 
  15 & COVIDhub-baseline & 2.00 & ba4\_ba5 & 19.69 & 21.67 & 0.97 & 1.00 & 1.00 & 1.00 \\ 
  16 & THieF\_12wk-noTransform & 2.00 & ba4\_ba5 & 21.52 & 32.82 & 0.68 & 1.00 & 1.09 & 1.51 \\ 
  17 & COVIDhub-4\_week\_ensemble & 3.00 & ba4\_ba5 & 15.22 & 22.86 & 0.76 & 0.99 & 0.64 & 0.83 \\ 
  18 & THieF\_6wk-4root & 3.00 & ba4\_ba5 & 16.82 & 23.30 & 0.72 & 0.99 & 0.70 & 0.85 \\ 
  19 & THieF\_ensemble-train3 & 3.00 & ba4\_ba5 & 17.95 & 25.29 & 0.76 & 1.00 & 0.75 & 0.92 \\ 
  20 & THieF\_ensemble-mean & 3.00 & ba4\_ba5 & 18.68 & 26.22 & 0.78 & 1.00 & 0.78 & 0.95 \\ 
  21 & THieF\_6wk-noTransform & 3.00 & ba4\_ba5 & 21.59 & 31.20 & 0.78 & 1.00 & 0.90 & 1.14 \\ 
  22 & sarima\_s7-noTransform & 3.00 & ba4\_ba5 & 22.06 & 31.34 & 0.78 & 0.99 & 0.92 & 1.14 \\ 
  23 & COVIDhub-baseline & 3.00 & ba4\_ba5 & 23.93 & 27.49 & 0.97 & 1.00 & 1.00 & 1.00 \\ 
  24 & THieF\_12wk-noTransform & 3.00 & ba4\_ba5 & 25.38 & 39.28 & 0.70 & 1.00 & 1.06 & 1.43 \\ 
  25 & COVIDhub-4\_week\_ensemble & 4.00 & ba4\_ba5 & 18.26 & 27.51 & 0.76 & 0.99 & 0.66 & 0.84 \\ 
  26 & THieF\_6wk-4root & 4.00 & ba4\_ba5 & 19.92 & 26.57 & 0.75 & 1.00 & 0.72 & 0.81 \\ 
  27 & THieF\_ensemble-train3 & 4.00 & ba4\_ba5 & 21.50 & 29.63 & 0.78 & 1.00 & 0.78 & 0.90 \\ 
  28 & THieF\_ensemble-mean & 4.00 & ba4\_ba5 & 22.13 & 30.41 & 0.79 & 1.00 & 0.81 & 0.93 \\ 
  29 & THieF\_6wk-noTransform & 4.00 & ba4\_ba5 & 24.17 & 34.33 & 0.80 & 1.00 & 0.88 & 1.05 \\ 
  30 & sarima\_s7-noTransform & 4.00 & ba4\_ba5 & 25.95 & 36.72 & 0.80 & 0.99 & 0.94 & 1.12 \\ 
  31 & COVIDhub-baseline & 4.00 & ba4\_ba5 & 27.46 & 32.80 & 0.97 & 1.00 & 1.00 & 1.00 \\ 
  32 & THieF\_12wk-noTransform & 4.00 & ba4\_ba5 & 29.06 & 45.41 & 0.69 & 1.00 & 1.06 & 1.38 \\ 
  33 & sarima\_s7-noTransform & 1.00 & delta & 10.38 & 14.69 & 0.70 & 0.97 & 0.30 & 0.30 \\ 
  34 & THieF\_ensemble-mean & 1.00 & delta & 11.60 & 18.76 & 0.64 & 0.99 & 0.34 & 0.39 \\ 
  35 & THieF\_ensemble-train3 & 1.00 & delta & 11.87 & 19.30 & 0.60 & 0.99 & 0.35 & 0.40 \\ 
  36 & THieF\_6wk-4root & 1.00 & delta & 14.16 & 22.68 & 0.51 & 0.95 & 0.42 & 0.47 \\ 
  37 & THieF\_12wk-noTransform & 1.00 & delta & 14.89 & 23.52 & 0.52 & 0.95 & 0.44 & 0.49 \\ 
  38 & THieF\_6wk-noTransform & 1.00 & delta & 15.57 & 25.83 & 0.51 & 0.97 & 0.46 & 0.54 \\ 
  39 & sarima\_s7-noTransform & 2.00 & delta & 14.64 & 20.80 & 0.69 & 0.98 & 0.27 & 0.27 \\ 
  40 & THieF\_ensemble-mean & 2.00 & delta & 16.33 & 26.27 & 0.65 & 0.99 & 0.30 & 0.34 \\ 
  41 & THieF\_ensemble-train3 & 2.00 & delta & 16.52 & 26.76 & 0.64 & 0.98 & 0.30 & 0.34 \\ 
  42 & THieF\_6wk-4root & 2.00 & delta & 17.99 & 29.88 & 0.54 & 0.98 & 0.33 & 0.38 \\ 
  43 & THieF\_12wk-noTransform & 2.00 & delta & 19.35 & 30.10 & 0.60 & 0.95 & 0.35 & 0.39 \\ 
  44 & THieF\_6wk-noTransform & 2.00 & delta & 19.58 & 32.20 & 0.62 & 0.97 & 0.36 & 0.41 \\ 
  45 & sarima\_s7-noTransform & 3.00 & delta & 23.20 & 35.71 & 0.60 & 0.97 & 0.32 & 0.34 \\ 
  46 & THieF\_ensemble-mean & 3.00 & delta & 25.61 & 41.49 & 0.58 & 0.97 & 0.35 & 0.40 \\ 
  47 & THieF\_ensemble-train3 & 3.00 & delta & 25.77 & 41.94 & 0.57 & 0.98 & 0.35 & 0.40 \\ 
  48 & THieF\_6wk-4root & 3.00 & delta & 27.03 & 44.49 & 0.52 & 0.99 & 0.37 & 0.43 \\ 
  49 & THieF\_12wk-noTransform & 3.00 & delta & 27.96 & 44.22 & 0.51 & 0.92 & 0.38 & 0.42 \\ 
  50 & THieF\_6wk-noTransform & 3.00 & delta & 28.40 & 46.47 & 0.53 & 0.94 & 0.39 & 0.44 \\ 
  51 & sarima\_s7-noTransform & 4.00 & delta & 30.50 & 49.25 & 0.54 & 0.94 & 0.35 & 0.39 \\ 
  52 & THieF\_ensemble-mean & 4.00 & delta & 31.71 & 51.12 & 0.54 & 0.97 & 0.36 & 0.40 \\ 
  53 & THieF\_ensemble-train3 & 4.00 & delta & 31.80 & 51.46 & 0.53 & 0.97 & 0.36 & 0.41 \\ 
  54 & THieF\_6wk-4root & 4.00 & delta & 32.28 & 52.56 & 0.52 & 0.99 & 0.37 & 0.41 \\ 
  55 & THieF\_6wk-noTransform & 4.00 & delta & 34.18 & 55.44 & 0.47 & 0.93 & 0.39 & 0.44 \\ 
  56 & THieF\_12wk-noTransform & 4.00 & delta & 34.38 & 54.48 & 0.49 & 0.90 & 0.39 & 0.43 \\ 
  57 & COVIDhub-4\_week\_ensemble & 1.00 & omicron & 25.02 & 37.57 & 0.57 & 0.92 & 0.74 & 0.78 \\ 
  58 & THieF\_ensemble-train3 & 1.00 & omicron & 31.71 & 46.62 & 0.40 & 0.86 & 2.19 & 2.81 \\ 
  59 & THieF\_ensemble-mean & 1.00 & omicron & 31.97 & 46.75 & 0.41 & 0.85 & 2.21 & 2.82 \\ 
  60 & COVIDhub-baseline & 1.00 & omicron & 34.00 & 48.18 & 0.66 & 0.94 & 1.00 & 1.00 \\ 
  61 & sarima\_s7-noTransform & 1.00 & omicron & 35.20 & 47.14 & 0.45 & 0.78 & 2.43 & 2.84 \\ 
  62 & THieF\_6wk-4root & 1.00 & omicron & 36.90 & 53.71 & 0.35 & 0.80 & 2.55 & 3.24 \\ 
  63 & THieF\_6wk-noTransform & 1.00 & omicron & 38.16 & 54.61 & 0.34 & 0.78 & 2.63 & 3.29 \\ 
  64 & THieF\_12wk-noTransform & 1.00 & omicron & 41.89 & 58.55 & 0.36 & 0.77 & 2.89 & 3.53 \\ 
  65 & COVIDhub-4\_week\_ensemble & 2.00 & omicron & 39.33 & 58.67 & 0.56 & 0.91 & 0.72 & 0.75 \\ 
  66 & THieF\_6wk-4root & 2.00 & omicron & 48.67 & 71.22 & 0.33 & 0.81 & 2.47 & 3.29 \\ 
  67 & THieF\_ensemble-train3 & 2.00 & omicron & 49.00 & 71.16 & 0.38 & 0.82 & 2.49 & 3.28 \\ 
  68 & THieF\_ensemble-mean & 2.00 & omicron & 49.04 & 70.92 & 0.39 & 0.82 & 2.49 & 3.27 \\ 
  69 & THieF\_6wk-noTransform & 2.00 & omicron & 53.99 & 77.42 & 0.33 & 0.76 & 2.74 & 3.57 \\ 
  70 & COVIDhub-baseline & 2.00 & omicron & 54.64 & 77.98 & 0.64 & 0.92 & 1.00 & 1.00 \\ 
  71 & THieF\_12wk-noTransform & 2.00 & omicron & 55.13 & 78.44 & 0.36 & 0.76 & 2.80 & 3.62 \\ 
  72 & sarima\_s7-noTransform & 2.00 & omicron & 59.62 & 78.63 & 0.44 & 0.74 & 3.03 & 3.63 \\ 
  73 & COVIDhub-4\_week\_ensemble & 3.00 & omicron & 54.85 & 81.62 & 0.53 & 0.89 & 0.76 & 0.78 \\ 
  74 & THieF\_6wk-4root & 3.00 & omicron & 62.25 & 90.82 & 0.32 & 0.80 & 2.60 & 3.30 \\ 
  75 & THieF\_ensemble-mean & 3.00 & omicron & 66.37 & 95.65 & 0.37 & 0.79 & 2.77 & 3.48 \\ 
  76 & THieF\_ensemble-train3 & 3.00 & omicron & 66.58 & 96.20 & 0.36 & 0.79 & 2.78 & 3.50 \\ 
  77 & THieF\_12wk-noTransform & 3.00 & omicron & 69.01 & 98.88 & 0.34 & 0.74 & 2.88 & 3.60 \\ 
  78 & THieF\_6wk-noTransform & 3.00 & omicron & 69.99 & 100.15 & 0.32 & 0.74 & 2.92 & 3.64 \\ 
  79 & COVIDhub-baseline & 3.00 & omicron & 72.69 & 104.38 & 0.60 & 0.90 & 1.00 & 1.00 \\ 
  80 & sarima\_s7-noTransform & 3.00 & omicron & 82.60 & 109.30 & 0.42 & 0.72 & 3.45 & 3.98 \\ 
  81 & COVIDhub-4\_week\_ensemble & 4.00 & omicron & 68.91 & 100.99 & 0.50 & 0.87 & 0.79 & 0.80 \\ 
  82 & THieF\_6wk-4root & 4.00 & omicron & 75.30 & 110.30 & 0.29 & 0.77 & 2.74 & 3.36 \\ 
  83 & THieF\_12wk-noTransform & 4.00 & omicron & 80.39 & 116.31 & 0.32 & 0.73 & 2.93 & 3.55 \\ 
  84 & THieF\_ensemble-mean & 4.00 & omicron & 80.41 & 116.55 & 0.35 & 0.77 & 2.93 & 3.55 \\ 
  85 & THieF\_ensemble-train3 & 4.00 & omicron & 80.87 & 117.46 & 0.34 & 0.77 & 2.94 & 3.58 \\ 
  86 & THieF\_6wk-noTransform & 4.00 & omicron & 83.75 & 120.76 & 0.30 & 0.72 & 3.05 & 3.68 \\ 
  87 & COVIDhub-baseline & 4.00 & omicron & 87.33 & 126.66 & 0.58 & 0.90 & 1.00 & 1.00 \\ 
  88 & sarima\_s7-noTransform & 4.00 & omicron & 102.60 & 136.89 & 0.41 & 0.70 & 3.74 & 4.17 \\ 
   \hline
\hline
\end{longtable}\caption{Summary of model performance stratified by horizon week and pandemic waves for the states level, aggregated across forecast week. The results from this table were used to create Figure 5 that summarized average WIS by horizon week and pandemic wave for the US national level and the states level. The table is ordered by ascending pandemic wave, horizon week, and WIS, with THieF-4wk-4root again having the lowest WIS for all horizons of the alpha wave and shorter horizons of the Winter 21 wave. The THieF-12wk-noTrans has the lowest WIS for 3 week ahead and 4 week ahead horizons during the Winter 21 wave. The baseline only has the best WIS for the 1 week ahead horizon during the Winter 21 wave.}
\end{center}
\end{table}



%%
%% Beginning of back matter
\backmatter  %% <--- mandatory

%%
%% We don't support endnotes

%%
%% A bibliography is required.
\interlinepenalty=10000  % prevent split bibliography entries
\bibliographystyle{umassthesis}
\bibliography{umthsmpl}
\end{document}

%%% Local Variables:
%%% mode: latex
%%% TeX-master: t
%%% End:

